\section{Milestone W4: the normal homological equation}\label{bpW:w4}

One step of the KAM iteration (milestone W5) removes from the perturbation its part of order at most two in the
normal variables by the time-one map of a Hamiltonian $F$ that solves the linearized equation
$\{F,N\}+\hat N=R$ of P96-K \S2. This part solves that equation for the two-frequency wave setting ($n=2$,
$d=1$, $\delta=-1$, $\bar p=p+1$) and proves the estimates of P96-K Lemmas~1 and~2 with explicit constants, including
the Lipschitz dependence on the parameter $\xi$.

\paragraph{Route.} The equation splits into seven components (scalar, tangential-linear, two normal-vector and three
normal-operator components). In every component the solution is a Fourier multiplier in the angle $\theta$: the
$k$-th Fourier coefficient is divided by a small divisor $\langle k,\omega\rangle+\langle l,\Omega\rangle$, entrywise
for the vector and operator components. One abstract lemma (Lemma~\ref{bpW:w4:solve}) handles all components. It
uses only the Banach-valued coefficient bound $\norm{\hat f_k}\le m\,e^{-2\pi\abs k_1s}$ and the explicit lattice
sums of Lemma~\ref{bpW:w4:lattice}; the operator-valued division is bounded by the Hilbert-matrix (Schur) lemma of
P96-K App.~A. P96-K p.~127 instead sums operator-valued coefficients by a square-sum bound which is false
(Remark~\ref{bpW:w4:counterexample}); the replacement keeps the power $\sigma^{-(2\tau+n+1)}$ of P96-K (4), so the
later smallness exponents of Theorem~A are unchanged (Remark~\ref{bpW:w4:compare}). The size of the components
of $R$ is read off from the weighted field norm $\abs{X_R}_r$ by evaluating the field at chosen points and
extracting polynomial coefficients by averages over roots of unity (Lemmas~\ref{bpW:w4:extract},
\ref{bpW:w4:components}); no Cauchy integral is needed. Derivatives of $F$ in $\theta$ are obtained by solving the
same equations with the $\theta$-derivatives of the components of $R$ as data, so no factor $\abs k$ is lost.

\paragraph{Inputs.} From W1: the weighted spaces and transports (Definition~\ref{bpW:w1:transport},
Lemma~\ref{bpW:w1:space}), the KAM phase space, domains and norms (Definition~\ref{bpW:w1:kam-space},
Remark~\ref{bpW:w1:kam-conventions}), the Hilbert-matrix bound (Lemma~\ref{bpW:w1:schur}), the spectral tails
(Lemma~\ref{bpW:w1:tails}), the sup-weighted sequences (Definition~\ref{bpW:w1:linfty}), and the Banach-valued
Fourier Lemma~\ref{bpW:w1:fourier-banach} with the scalar uniqueness of Lemma~\ref{bpW:w1:fourier-l2}(iii) (the exact
content used is listed in Remark~\ref{bpW:w4:w1-requests}). From W0: the action coefficients
(Theorem~\ref{bpW:w0n:AB}). From W3: the form of the Hamiltonian (Remark~\ref{bpW:w3:interface}); see
\S\ref{bpW:w4:interface-sec}.

\paragraph{Regularity used.} W4 uses only complex differentiability (\Lean{DifferentiableOn Complex}) of the
$\theta$-dependent components of $R$ on a strip, and boundedness. It does not use analyticity of $X_R$ on an
infinite-dimensional domain; P96-K Lemma~1 assumes ``$X_R$ real analytic'', which is more than is used here. In
particular W4 does not depend on the pending upgrade of the W3 interface from holomorphic to analytic fields
(review finding W3-R1); that dependency, if any, sits in W5, where $R$ is produced from $P$.

\subsection{External facts}\label{bpW:w4:ext}

Besides the external facts of W1 (\S\ref{bpW:w1:ext}), this part uses:
{\raggedright
\begin{enumerate}[label=(W4-E\arabic*),ref=(W4-E\arabic*)]
\item\label{bpW:w4:ext-roots} \emph{Roots of unity.} For $m\in\Z$, $\sum_{u=0}^{3}i^{um}=4$ if $4\mid m$ and $=0$
  otherwise (finite geometric sum; Mathlib \Lean{Finset.sum\_range} with \Lean{geom\_sum\_eq}).
\item\label{bpW:w4:ext-parallelogram} \emph{Parallelogram law} in a complex inner product space:
  $\norm{x+y}^2+\norm{x-y}^2=2\norm x^2+2\norm y^2$ (\Lean{parallelogram\_law\_with\_norm}).
\item\label{bpW:w4:ext-exp} \emph{Exponential inequalities.} $e^t\ge1+t$ for real $t$
  (\Lean{Real.add\_one\_le\_exp}); hence $e^{t-1}\ge t$, and for $q>0$, $x\ge0$: $x^qe^{-x}\le(q/e)^q$. Also
  $\sum_{N\ge1}e^{-xN}=1/(e^x-1)$ for $x>0$ (\Lean{tsum\_geometric\_of\_lt\_one}).
\item\label{bpW:w4:ext-intparts} \emph{Periodic integration by parts.} If $g:\R\to E$ ($E$ Banach) is $C^1$ and
  $1$-periodic, then $\int_0^1g'=0$ (\Lean{intervalIntegral.integral\_eq\_sub\_of\_hasDerivAt}); continuous linear maps
  commute with Bochner integrals (\Lean{ContinuousLinearMap.integral\_comp\_comm}).
\item\label{bpW:w4:ext-riesz} \emph{Riesz representation of bounded sesquilinear forms.} Let $E$ be a complex
  Hilbert space (Mathlib: \Lean{[NormedAddCommGroup E] [InnerProductSpace Complex E] [CompleteSpace E]}; here
  $E=\ell^2_\C$, complete by \ref{bpW:w1:ext-lp}) and $B:E\to E\to\C$ a bounded form, conjugate-linear in the first
  and linear in the second argument (Mathlib type \Lean{E \ensuremath{\to}L\ensuremath{\star}[Complex] E \ensuremath{\to}L[Complex] Complex}). Then
  $B^\sharp=$ \Lean{InnerProductSpace.continuousLinearMapOfBilin B} is a continuous linear map $E\to E$ with
  $\langle B^\sharp v,w\rangle=B(v,w)$ for all $v,w$ (\Lean{InnerProductSpace.continuousLinearMapOfBilin\_apply}), where
  $\langle\cdot,\cdot\rangle$ is the inner product, conjugate-linear in the first slot. Moreover
  $\norm u=\sup_{\norm w\le1}\abs{\langle u,w\rangle}$ (\Lean{InnerProductSpace.toDual} is a linear isometry equivalence
  $E\simeq E^*$, and $\langle u,w\rangle=$ \Lean{toDual Complex E u w}, \Lean{InnerProductSpace.toDual\_apply\_apply}).
  (Checked in Mathlib, \texttt{Mathlib/Analysis/InnerProductSpace/Dual.lean}.)
\end{enumerate}
\par}
The bound in~\ref{bpW:w4:ext-exp} follows from $e^{x/q-1}\ge x/q$ raised to the power $q$.

\subsection{Setting and hypotheses}\label{bpW:w4:setting}

\begin{convention}[Normal modes, reindexing, coordinates]\label{bpW:w4:index}
Fix $a\ge0$, $p\ge0$ and put $\bar p=p+1$ (the wave case of P96-K Assumption~C, $\bar p-p=1$). All results of this part hold for $a,p\ge0$
except the identification of the weak bracket with a derivative of $N$, which needs the guard \eqref{bpW:w4:guard}
($a>0$ or $p\ge\frac12$; Lemma~\ref{bpW:w4:lie}(iv)); the wave consumer $a=p=1$ satisfies it. The tangential modes
are the wave modes $1,2$; the normal modes are the wave modes $J\ge3$. Following P96-K (index $j\ge1$) we write
\[J=j+2\qquad(j\ge1)\]
and index every normal quantity by the \emph{P96-K index} $j$: $\Omega_j$, $z_j$, $\bar z_j$, matrix entries
$M_{ij}$. For $q\in\R$ let $\ell^N_{a,q}$ be the space of complex sequences $(z_j)_{j\ge1}$ with
\[\norm z_{a,q}^2=\sum_{j\ge1}\omega_{a,q}(j+2)^2\abs{z_j}^2<\infty,\qquad\omega_{a,q}(J)=J^qe^{aJ}.\]
It is realized in Lean as the closed subspace of \Lean{lp (fun \_ : PNat => Complex) 2} (weighted coordinates at
weight $(a,q)$, Definition~\ref{bpW:w1:transport}) of vectors whose entries at $J=1,2$ vanish, the entry at $J$
being the coordinate $j=J-2$; this is the convention of Lemma~\ref{bpW:w1:tails}. The weights are the wave weights
$\omega_{a,q}(J)$ at the wave index, not P96-K's $j^qe^{aj}$; see Remark~\ref{bpW:w4:weights}. We write
$\norm\cdot=\norm\cdot_{a,p}$ on the normal variables and $\norm\cdot_{\bar p}=\norm\cdot_{a,\bar p}$ on gradients.
On $\C^2$ (angles $\theta$, actions $y$, frequencies $\omega$) $\abs\cdot$ is the max norm; on $\Z^2$,
$\abs k_1=\abs{k_1}+\abs{k_2}$.

The unweighted bilinear pairing is $\langle x,x'\rangle_0=\sum_{j\ge1}x_jx'_j$ (no conjugation), and on $\C^2$,
$\langle b,y\rangle=b_1y_1+b_2y_2$. Since $\omega_{a,q}(J)\,\omega_{a,q'}(J)=J^{q+q'}e^{2aJ}\ge1$ whenever
$q+q'\ge0$, Cauchy--Schwarz gives
\begin{equation}\label{bpW:w4:pairing}
\abs{\langle x,x'\rangle_0}\le\norm x_{a,q}\norm{x'}_{a,q'}\qquad(q+q'\ge0).
\end{equation}
The normal phase variables are $(u,v)\in\ell^N_{a,p}\times\ell^N_{a,p}$ (P96-K's $u,v$; wave $u_j=q_{j+2}$,
$v_j=p_{j+2}$), and
\[z=\frac{u+iv}{\sqrt2},\qquad\bar z=\frac{u-iv}{\sqrt2},\qquad u=\frac{z+\bar z}{\sqrt2},\qquad v=\frac{z-\bar z}{i\sqrt2},\]
with $z,\bar z$ independent complex variables on complex domains. This is the W0 convention
(Definition~\ref{bpW:w0:z}); P96-K p.~124 uses $z=(u-iv)/\sqrt2$, which exchanges $z$ and $\bar z$
(Remark~\ref{bpW:w1:kam-conventions}). From these formulas,
\begin{equation}\label{bpW:w4:zu}
\norm z,\norm{\bar z}\le\frac{\norm u+\norm v}{\sqrt2},\qquad\norm u+\norm v\le\sqrt2\,(\norm z+\norm{\bar z}).
\end{equation}
The angle $\theta\in\C^2$ has period one ($x=2\pi\theta$ is P96-K's angle); the strip is
$T_s=\{\theta\in\C^2:\abs{\operatorname{Im}\theta}<s\}$ with $s=s_K/(2\pi)$, and the domain $D(s,r)$ is that of
Definition~\ref{bpW:w1:kam-space} (with the normal norms above). The symplectic form is the pull-back
$2\pi\sum_id\theta_i\wedge dy_i+\sum_jdu_j\wedge dv_j$, so Hamilton's equations are
\begin{equation}\label{bpW:w4:hameq}
\dot\theta=\frac{1}{2\pi}H_y,\quad\dot y=-\frac{1}{2\pi}H_\theta,\quad\dot u=H_v,\quad\dot v=-H_u,\qquad
\text{equivalently}\quad\dot z=-iH_{\bar z},\quad\dot{\bar z}=iH_z,
\end{equation}
where $H_u,H_v,H_z,H_{\bar z}$ are gradients for $\langle\cdot,\cdot\rangle_0$ and
$H_z=(H_u-iH_v)/\sqrt2$, $H_{\bar z}=(H_u+iH_v)/\sqrt2$, $H_u=(H_z+H_{\bar z})/\sqrt2$, $H_v=i(H_z-H_{\bar z})/\sqrt2$.
Frequencies are angular: $\omega$ is P96-K's $\omega$, and the period-one frequency of the frozen statement is
$\nu=\omega/(2\pi)$.
\end{convention}
\source{P96-K pp.~120, 124 (phase space $\mathcal P^{a,p}$, $z=(u-iv)/\sqrt2$); conversions ours (W1
Definition~\ref{bpW:w1:kam-space}, Remark~\ref{bpW:w1:kam-conventions}). Lean guards: no division; $\theta$-periodicity
is \Lean{Function.Periodic} in each coordinate.}

\begin{remark}[Why the wave weights are kept]\label{bpW:w4:weights}
P96-K weights the $j$-th normal coordinate by $j^pe^{aj}$; the wave space weights it by $(j+2)^pe^{a(j+2)}$. The
two norms are equivalent within the factor $3^pe^{2a}$ (Remark~\ref{bpW:w3:interface}). Every estimate of this part
is insensitive to the choice of positive weights (Lemma~\ref{bpW:w4:schur}) except one: the diagonal of an operator
$\ell^N_{a,p}\to\ell^N_{a,\bar p}$ is a multiplier $\mu_j$ with $\abs{\mu_j}\,(j+2)^{\bar p-p}\le\norm{\cdot}$, and
$(j+2)\ge j$ is all that P96-K needs there (Theorem~\ref{bpW:w4:lemma1}(iv)). So W4 works with the wave weights
throughout and never incurs the factor $3^pe^{2a}$.
\end{remark}

\begin{definition}[Normal operators and their norms]\label{bpW:w4:ops}
$\mathcal L$ is the Banach space of bounded linear maps $\ell^N_{a,p}\to\ell^N_{a,\bar p}$ with the operator norm
$\norm\cdot$. For $M\in\mathcal L$ its matrix is $M_{ij}=\langle Me_j,e_i\rangle_0$, where $e_j$ is the $j$-th unit
sequence; $M^T$ is the map with matrix $(M_{ji})$ when it lies in $\mathcal L$; $M$ is \emph{symmetric} if $M^T=M$.
The \emph{bilinear norm} of a bounded bilinear form $\beta$ on $\ell^N_{a,p}$ is
$\norm\beta_{\rm bil}=\sup\{\abs{\beta(z,z')}:\norm z=\norm{z'}=1\}$; $\mathcal B$ is the Banach space of these forms,
and $\Lambda:\mathcal L\to\mathcal B$, $\Lambda M(z,z')=\langle Mz,z'\rangle_0$, satisfies
$\norm{\Lambda M}_{\rm bil}\le\norm M$ by \eqref{bpW:w4:pairing}. The \emph{dual norm} of $c$ is
$\norm c_*=\sup\{\abs{\langle c,z\rangle_0}:\norm z=1\}=\bigl(\sum_j\abs{c_j}^2\omega_{a,p}(j+2)^{-2}\bigr)^{1/2}$, and
$\ell^N_*$ is the space of $c$ with $\norm c_*<\infty$ (a weighted $\ell^2$, Lean: \Lean{lp} with weight
$\omega_{a,p}(j+2)^{-1}$).
\end{definition}
\source{ours (P96-K p.~127 speaks of ``bounded linear operators $\ell^{a,p}\to\ell^{a,\bar p}$'' and their
``operator norm $\norm\cdot$'').}

\begin{definition}[Taylor Hamiltonians]\label{bpW:w4:taylor}
Let $s>0$. The class $\mathcal T_s$ consists of the functions
\begin{equation}\label{bpW:w4:R}
R(\theta,y,z,\bar z)=a(\theta)+\langle b(\theta),y\rangle+\langle c^+(\theta),z\rangle_0+\langle c^-(\theta),\bar z\rangle_0
+\langle P(\theta)z,z\rangle_0+\langle Q(\theta)z,\bar z\rangle_0+\langle P'(\theta)\bar z,\bar z\rangle_0
\end{equation}
whose \emph{components} $a:T_s\to\C$, $b:T_s\to\C^2$, $c^\pm:T_s\to\ell^N_{a,\bar p}$, $P,P',Q,Q^T:T_s\to\mathcal L$ are
complex differentiable, $1$-periodic in each $\theta_i$, and bounded, with $P(\theta)$, $P'(\theta)$ symmetric. For
$R\in\mathcal T_s$ and $(\theta,y,z,\bar z)\in T_s\times\C^2\times\ell^N_{a,p}\times\ell^N_{a,p}$ put
\[R_y=b,\quad R_\theta=(\partial_{\theta_1}R,\partial_{\theta_2}R),\quad R_z=c^++2Pz+Q^T\bar z,\quad
R_{\bar z}=c^-+Qz+2P'\bar z,\]
with $R_u,R_v$ as in Convention~\ref{bpW:w4:index}; $R_z,R_{\bar z},R_u,R_v\in\ell^N_{a,\bar p}$. The \emph{field}
is $X_R=\bigl(\frac1{2\pi}R_y,-\frac1{2\pi}R_\theta,R_v,-R_u\bigr)$ and, for $r>0$,
\[\abs{X_R}_r=\frac1{2\pi}\abs{R_y}+\frac{1}{2\pi r^2}\abs{R_\theta}+\frac1r\bigl(\norm{R_u}_{\bar p}+\norm{R_v}_{\bar p}\bigr),\]
the norm of Definition~\ref{bpW:w1:kam-space} in period-one coordinates. $R$ is \emph{real} if for real $\theta$:
$a(\theta)\in\R$, $b(\theta)\in\R^2$, $c^-(\theta)=\overline{c^+(\theta)}$, $P'(\theta)=\overline{P(\theta)}$ and
$Q(\theta)^T=\overline{Q(\theta)}$ (entrywise conjugation); equivalently $R$ is real at real points
$\theta\in\R^2$, $y\in\R^2$, $\bar z=\overline z$. The \emph{mean} of a component $g$ is
$[g]=\hat g_0=\int_{[0,1]^2}g(x)\,dx$ (Bochner integral over the real torus).
\end{definition}
\source{P96-K (3) p.~124 (``$R=\sum_{2\abs m+\abs{q+\bar q}\le2}\sum_kR_{kmq\bar q}e^{i\langle k,x\rangle}y^mz^q\bar z^{\bar q}$'')
and the decomposition $R=R^0+R^1+R^2$ on p.~125; here $R^{00}=a+\langle b,y\rangle$, $R^{10}=c^+$, $R^{01}=c^-$,
$R^{20}=P$, $R^{11}=Q$, $R^{02}=P'$. The derivatives are the gradients of the polynomial \eqref{bpW:w4:R}: for
symmetric $P$, $D_z\langle Pz,z\rangle_0h=2\langle Pz,h\rangle_0$, and $D_z\langle Qz,\bar z\rangle_0h=\langle Q^T\bar z,h\rangle_0$.}

\begin{definition}[Frequencies, divisors, Melnikov conditions]\label{bpW:w4:melnikov}
Let $\Pi\subset\R^2$, and let $\omega:\Pi\to\R^2$, $\Omega=(\Omega_j)_{j\ge1}:\Pi\to\R^{\N_+}$ satisfy, for constants
$M_0,M_\omega,M_\Omega\ge0$ and all $\xi,\zeta\in\Pi$,
\begin{equation}\label{bpW:w4:freq}
\Omega_j=j+2+\tilde\Omega_j,\quad\sup_jj\abs{\tilde\Omega_j(\xi)}\le M_0,\quad\abs{\omega(\xi)-\omega(\zeta)}\le M_\omega\abs{\xi-\zeta},\quad
\sup_jj\abs{\Omega_j(\xi)-\Omega_j(\zeta)}\le M_\Omega\abs{\xi-\zeta},
\end{equation}
and put $M=M_\omega+M_\Omega$ (P96-K's $M=\abs\omega^{\mathcal L}+\abs\Omega^{\mathcal L}_{-\delta}$, $-\delta=1$). Let
$\tau\ge0$, $A_0=1$ and $A_k=1+\abs k_1^\tau$ for $k\ne0$ (so $A_0=1$ also when $\tau=0$; this is the
convention throughout this part), $\abs l=\sum_j\abs{l_j}$ and $\jap l=\max(1,\abs{\sum_jjl_j})$ for finitely
supported $l\in\Z^{\N_+}$, and
\[\mathcal Z=\{(k,l)\in\Z^2\times\Z^{\N_+}:(k,l)\ne0,\ \abs l\le2\}.\]
For $\alpha>0$, the \emph{Melnikov conditions} $(\mathrm H_\alpha)$ are
\[\abs{\langle k,\omega(\xi)\rangle+\langle l,\Omega(\xi)\rangle}\ge\alpha\frac{\jap l}{A_k}\qquad\text{for all }\xi\in\Pi,\ (k,l)\in\mathcal Z.\]
The divisors used below are, for $k\in\Z^2$ and $i,j\ge1$ (with $\xi$ suppressed):
\[\begin{array}{lll}
\delta^0_k=\langle k,\omega\rangle\ (k\ne0),&\delta^+_{k,j}=\langle k,\omega\rangle-\Omega_j,&\delta^-_{k,j}=\langle k,\omega\rangle+\Omega_j,\\[2pt]
\delta^P_{k,ij}=\langle k,\omega\rangle-\Omega_i-\Omega_j,&\delta^{P'}_{k,ij}=\langle k,\omega\rangle+\Omega_i+\Omega_j,&
\delta^Q_{k,ij}=\langle k,\omega\rangle+\Omega_i-\Omega_j\ ((k,i)\ne(0,j)),
\end{array}\]
with the \emph{weights} $w^0=1$, $w^\pm_j=j$, $w^P_{ij}=w^{P'}_{ij}=i+j$, $w^Q_{ij}=\max(1,\abs{i-j})$.
\end{definition}
\source{P96-K p.~121 (``$\jap l_d=\max(1,\abs{\sum j^dl_j})$, $A_k=1+\abs k^\tau$'', $Z$), Lemma~1 p.~125 (the
hypothesis), Assumption~B p.~120; $d=1$. The divisor list is P96-K p.~125 (``$\langle k,\omega\rangle+\langle q-\bar q,\Omega\rangle$'')
in the W0 sign convention, see Lemma~\ref{bpW:w4:lie}(iii). P96-K's $\abs k$ for $k\in\Z^n$ is not specified; we use
$\abs k_1$, which is what the lattice sums of Lemma~\ref{bpW:w4:lattice} need.}

\begin{lemma}[The wave frequencies]\label{bpW:w4:wave}
Let $\lambda_J=\sqrt{J^2+m}$, $b=(\lambda_1^{-1},\lambda_2^{-1})$, and for $\xi\in\R^2$
\[\omega(\xi)=(\lambda_1,\lambda_2)+A\xi,\qquad\Omega_j(\xi)=\lambda_{j+2}+\frac{3}{2\pi\lambda_{j+2}}\langle b,\xi\rangle\quad(j\ge1),\]
with $A$ of Theorem~\ref{bpW:w0n:AB}. Then for every $\Pi\subset\R^2$:
\begin{enumerate}[label=(\roman*)]
\item \eqref{bpW:w4:freq} holds with $M_\omega=\frac{9}{8\pi\lambda_1^2}+\frac{3}{2\pi\lambda_1\lambda_2}$,
  $M_\Omega=\frac3{2\pi}(\lambda_1^{-1}+\lambda_2^{-1})$ and $M_0=\frac m2+\frac3{2\pi}\sup_\Pi\abs{\langle b,\xi\rangle}$
  (when finite). In particular P96-K Assumption~B holds with $d=1$, $\delta=-1$ and the fixed sequence
  $\bar\Omega_j=j+2$, whose constant $2$ is a lower-order term allowed by Assumption~B.
\item (Reindexing.) Write $\jap l^{\rm w}=\max(1,\abs{\sum_j(j+2)l_j})$ for the wave index. For $1\le\abs l\le2$,
  $\jap l\le\jap l^{\rm w}\le3\jap l$ unless $\sum_jl_j=0$, in which case $\jap l^{\rm w}=\jap l$. Hence
  $(\mathrm H_\alpha)$ stated with $\jap l^{\rm w}$ implies $(\mathrm H_\alpha)$, and $(\mathrm H_{\alpha})$ implies the
  conditions with $\jap l^{\rm w}$ and $\alpha/3$.
\end{enumerate}
\end{lemma}

\begin{proof}
\proofstep{Lipschitz constants}\label{stepW:w4:waveLip}
$\abs{A(\xi-\zeta)}\le\max_i\sum_{i'}A_{ii'}\abs{\xi-\zeta}$, the maximal row sum being the first since
$\lambda_1<\lambda_2$ and all $A_{ii'}>0$ (Theorem~\ref{bpW:w0n:AB}). For $\Omega$,
$j\abs{\Omega_j(\xi)-\Omega_j(\zeta)}=\frac{3j}{2\pi\lambda_{j+2}}\abs{\langle b,\xi-\zeta\rangle}\le\frac3{2\pi}\abs b_1\abs{\xi-\zeta}$
since $\lambda_{j+2}>j$.
\proofstep{Tails}\label{stepW:w4:waveTail}
$\tilde\Omega_j=(\lambda_{j+2}-(j+2))+\frac{3\langle b,\xi\rangle}{2\pi\lambda_{j+2}}$ and
$0<\lambda_J-J=m/(\lambda_J+J)\le m/(2J)$, so $j\abs{\tilde\Omega_j}\le\frac m2+\frac3{2\pi}\abs{\langle b,\xi\rangle}$.
This is Lemma~\ref{bpW:w1:tails} in the P96-K index.
\proofstep{Index comparison}\label{stepW:w4:waveIdx}
$\sum_j(j+2)l_j=\sum_jjl_j+2\sum_jl_j$. If $\sum l_j=0$ the two agree. Otherwise, up to a global sign,
$l=e_j$ ($\sum jl_j=j$, $\sum(j+2)l_j=j+2\le3j$) or $l=e_i+e_j$ ($i+j$ and $i+j+4\le3(i+j)$), and the ratio lies in
$[1,3]$.
\end{proof}
\source{ours; P96-W p.~290 (frequencies of the normal form) with the corrected $A$, $B$ (E-PDE-W-1,
Theorem~\ref{bpW:w0n:AB}); P96-K Assumption~B p.~120. The source applies P96-K with normal index $j\ge1$ without
spelling out the shift $j\mapsto j+2$.}

\begin{lemma}[The conditions with $k=0$ for the wave frequencies]\label{bpW:w4:k0}
In the setting of Lemma~\ref{bpW:w4:wave}, let $\alpha_0(m)=2/\sqrt{9+m}$. If $\xi\in\R^2$ satisfies
$\abs{\langle b,\xi\rangle}\le2\pi$ and $0<\alpha\le\alpha_0(m)$, then $\abs{\langle l,\Omega(\xi)\rangle}\ge\alpha\jap l$ for
every $l$ with $1\le\abs l\le2$. That is, the conditions $(\mathrm H_\alpha)$ with $k=0$ hold on
$\Pi\subset\{\abs{\langle b,\xi\rangle}\le2\pi\}$.
\end{lemma}

\begin{proof}
Write $\beta=\frac3{2\pi}\langle b,\xi\rangle$, so $\abs\beta\le3$, and $J=j+2\ge3$, $I=i+2$.
\proofstep{One mode}\label{stepW:w4:k0one}
$\Omega_j\ge\lambda_J-\abs\beta/\lambda_J\ge J-1=j+1>\alpha\jap{e_j}$, since $\lambda_J>J\ge3$ and $\alpha\le1$. The
same bound gives $\Omega_i+\Omega_j\ge i+j+2>\alpha\jap{e_i+e_j}$.
\proofstep{Differences}\label{stepW:w4:k0diff}
For $i>j$: $\Omega_i-\Omega_j=(\lambda_I-\lambda_J)\bigl(1-\beta/(\lambda_I\lambda_J)\bigr)\ge\frac23(\lambda_I-\lambda_J)$, using
$\frac1{\lambda_I}-\frac1{\lambda_J}=-\frac{\lambda_I-\lambda_J}{\lambda_I\lambda_J}$ and $\lambda_I\lambda_J>9$. Further
$\lambda_I-\lambda_J=\frac{(I-J)(I+J)}{\lambda_I+\lambda_J}\ge\frac{I-J}{\sqrt{1+m/9}}$ because $\lambda_K\le K\sqrt{1+m/9}$ for
$K\ge3$. Hence $\Omega_i-\Omega_j\ge\alpha_0(m)(i-j)$, and $\jap{e_i-e_j}=i-j$. Signs of $l$ are symmetric.
\end{proof}
\source{ours. Superseded by Lemma~\ref{bpW:w7:sep} (W7) and used by no later result; kept for reference.
A numerical check over $j\le150$ and $m\in\{10^{-3},\frac12,1,10,10^3\}$ gives minimal ratio
$\ge1.12$ (\texttt{.drafts/W/bp-scratch-w4/checks.py}). This settles the sets $\mathcal R^0_{0l}(\alpha)$ of P96-K
Theorem~A for W7: they are empty for $\alpha\le\alpha_0(m)$.}

\begin{remark}[What is checked about the Melnikov conditions]\label{bpW:w4:melnikov-check}
The conditions with $k\ne0$ cannot hold for all $\xi$; P96-K Theorem~A removes the resonance zones and Theorem~D
(milestone W7) estimates their measure. What W4 checks is that the conditions of P96-K Lemma~1 are the right ones in
the wave setting: (a) they are stated in the P96-K index $j$ with $d=1$, and the reindexing changes $\jap l$ by at most
the factor $3$ (Lemma~\ref{bpW:w4:wave}(ii)); (b) the exchange $z\leftrightarrow\bar z$ between W0 and P96-K replaces
$l$ by $-l$, which maps $\mathcal Z$ to itself and preserves $\jap l$, so the conditions are the same; (c) every divisor
of Definition~\ref{bpW:w4:melnikov} is of the form $\langle k,\omega\rangle+\langle l,\Omega\rangle$ with
$(k,l)\in\mathcal Z$, and $\jap l\ge w$ (Lemma~\ref{bpW:w4:divisors}); (d) the conditions with $k=0$ hold outright for
small $\alpha$ (Lemma~\ref{bpW:w4:k0}).
\end{remark}

\subsection{Fourier series with Banach-space coefficients}\label{bpW:w4:fourier-sec}

\begin{remark}[The operator-valued square-sum bound of P96-K p.~127 is false]\label{bpW:w4:counterexample}
P96-K p.~127 states, for the operator-valued coefficients $\tilde R_k$ of $\tilde R=\partial_{\bar z}\partial_zR$,
``we have, as before, $\sum_k\abs{\!\abs{\!\abs{\tilde R_k}\!}\!}^2e^{2\abs k\abs s}\le2^n\abs{\!\abs{\!\abs{\tilde R}\!}\!}_{D(s)}$''
(``as before'' refers to the $L^2$ bound for $\ell^{a,\bar p}$-valued coefficients on p.~126, printed there with
$\norm{\dot R}^2_{D(s)}$; on p.~127 the square on the right side is missing). This is false, even with the square. In period-one notation, fix $s>0$, $N\ge1$ and let
$A_N(\theta)$ be the diagonal operator on $\ell^2$ with
\[A_N(\theta)e_j=e^{-2\pi js}e^{2\pi ij\theta_1}e_j\quad(1\le j\le N),\qquad A_N(\theta)e_j=0\quad(j>N).\]
It is entire and $1$-periodic, and on $T_s$ its entries have modulus $e^{-2\pi j(s+\operatorname{Im}\theta_1)}<1$, so
$\sup_{T_s}\norm{A_N(\theta)}=1$ (approached as $\operatorname{Im}\theta_1\to-s$). Its nonzero Fourier coefficients
are $\hat A_{(j,0)}=e^{-2\pi js}e_j\otimes e_j$, of norm $e^{-2\pi js}$, so
\[\sum_k\norm{\hat A_k}^2e^{4\pi\abs k_1s}=\sum_{j=1}^N1=N,\]
which exceeds $2^n\sup\norm{A_N}^2=4$ for $N\ge5$ ($n=2$; verified in
\texttt{.drafts/W/bp-scratch-w4/checks.py}). Diagonal operators between weighted spaces are transported to this
example by the diagonal maps of Definition~\ref{bpW:w1:transport}, and the example is the mixed Hessian of the quadratic
Hamiltonian $\sum_{j\le N}\mu_j(\theta)z_j\bar z_j$, so it lies in the class P96-K uses. The Hilbert-space bound
(Lemma~\ref{bpW:w1:fourier-l2}(i)) is true for $\mathcal H$-valued maps; it fails here because $\mathcal L$ with the
operator norm is not a Hilbert space. The valid replacement is the Banach-valued bound
$\norm{\hat A_k}\le\sup_{T_s}\norm A\,e^{-2\pi\abs k_1s}$ (attained here for every $k=(j,0)$), followed by an
$\ell^1$ summation after a strip loss. This is review finding W1-R1, rechecked here.
\end{remark}

\begin{remark}[What W4 uses from W1]\label{bpW:w4:w1-requests}
W4 uses exactly the following, for a complex Banach space $E$, $s>0$ and
$f:T_s\to E$ complex differentiable, $1$-periodic, with $\norm{f(\theta)}\le m$ on $T_s$, and with
$\hat f_k=\int_{[0,1]^2}f(x)e^{-2\pi i\langle k,x\rangle}dx$:
\begin{enumerate}[label=(F\arabic*),ref=(F\arabic*)]
\item\label{bpW:w4:F1} $\norm{\hat f_k}\le m\,e^{-2\pi\abs k_1s}$ for all $k\in\Z^2$ (contour shift with the sign
  correction of W1-R3);
\item\label{bpW:w4:F2} if $c_k\in E$ and $\sum_k\norm{c_k}e^{2\pi\abs k_1s'}<\infty$, then
  $g(\theta)=\sum_kc_ke^{2\pi i\langle k,\theta\rangle}$ converges absolutely and uniformly on $T_{s'}$, $g$ is complex
  differentiable and $1$-periodic there, $\hat g_k=c_k$, and $\norm{g(\theta)}\le\sum_k\norm{c_k}e^{2\pi\abs k_1s'}$;
\item\label{bpW:w4:F3} (uniqueness, scalar case $E=\C$) if $f$ is as above and $\hat f_k=0$ for all $k$, then $f=0$ on
  $T_{s'}$ for every $s'<s$, hence on $T_s$.
\end{enumerate}
\ref{bpW:w4:F1} is Lemma~\ref{bpW:w1:fourier-banach}(i), \ref{bpW:w4:F2} is Lemma~\ref{bpW:w1:fourier-banach}(iii), and
\ref{bpW:w4:F3} is Lemma~\ref{bpW:w1:fourier-l2}(iii) with $\mathcal H=\C$ (a Hilbert space; only scalar uniqueness is
used). The Hilbert-valued square-sum bound Lemma~\ref{bpW:w1:fourier-l2}(i) is not used for operators. The
$n=2$ lattice sums are proved below with sharper constants (Lemma~\ref{bpW:w4:lattice}); Lemma~\ref{bpW:w1:lattice}
gives the same powers of $\sigma$ for general $n$.
\end{remark}

\begin{lemma}[Coefficients of linear images and derivatives]\label{bpW:w4:fourier-ops}
Let $f:T_s\to E$ be as in Remark~\ref{bpW:w4:w1-requests}.
\begin{enumerate}[label=(\roman*)]
\item For a continuous linear $\Lambda:E\to E'$, $\widehat{(\Lambda f)}_k=\Lambda\hat f_k$.
\item For $i\in\{1,2\}$, $\partial_{\theta_i}f$ is complex differentiable and $1$-periodic on $T_s$, and
  $\widehat{(\partial_{\theta_i}f)}_k=2\pi ik_i\hat f_k$.
\item If $f$ depends on $\xi\in\Pi$ and $\norm{f(\theta,\xi)-f(\theta,\zeta)}\le m^{\mathcal L}\abs{\xi-\zeta}$ on $T_s$,
  then $\norm{\hat f_k(\xi)-\hat f_k(\zeta)}\le m^{\mathcal L}\abs{\xi-\zeta}e^{-2\pi\abs k_1s}$.
\end{enumerate}
\end{lemma}

\begin{proof}
\proofstep{Linear images}\label{stepW:w4:linimg}
\ref{bpW:w4:ext-intparts}: $\Lambda$ commutes with the integral defining $\hat f_k$.
\proofstep{Derivatives}\label{stepW:w4:deriv}
The derivative of a complex differentiable map is complex differentiable (\ref{bpW:w1:ext-onevar} in the variable
$\theta_i$; on the real torus it suffices that $f$ is $C^1$), and periodicity passes to $\partial_{\theta_i}f$. On the
real torus, $\partial_{x_i}\bigl(f(x)e^{-2\pi i\langle k,x\rangle}\bigr)=(\partial_{x_i}f)e^{-2\pi i\langle k,x\rangle}-2\pi ik_if(x)e^{-2\pi i\langle k,x\rangle}$
and the left side integrates to zero over a period in $x_i$ (\ref{bpW:w4:ext-intparts}); integrate over the other
variable (Fubini for continuous functions on $[0,1]^2$).
\proofstep{Parameters}\label{stepW:w4:fpar}
Apply \ref{bpW:w4:F1} to $f(\cdot,\xi)-f(\cdot,\zeta)$, whose coefficients are $\hat f_k(\xi)-\hat f_k(\zeta)$.
\end{proof}
\source{standard.}

\begin{lemma}[Lattice sums in $\Z^2$]\label{bpW:w4:lattice}
Let $\sigma>0$ and $S_q(\sigma)=\sum_{k\in\Z^2}\abs k_1^qe^{-2\pi\abs k_1\sigma}$ for $q\ge0$ (with $0^0=1$).
\begin{enumerate}[label=(\roman*)]
\item $\#\{k\in\Z^2:\abs k_1=N\}=4N$ for $N\ge1$.
\item $S_0(\sigma)=\coth(\pi\sigma)^2\le(1+(\pi\sigma)^{-1})^2$.
\item For $q>0$, $S_q(\sigma)\le\dfrac4{\pi\sigma}\Bigl(\dfrac{q+1}{e\pi\sigma}\Bigr)^{q+1}$.
\item Let $\tau\ge0$, $A_0=1$ and $A_k=1+\abs k_1^\tau$ for $k\ne0$ (as in Definition~\ref{bpW:w4:melnikov}), and
  \[B^{\rm W}_\sigma=\sum_{k\in\Z^2}(1+\abs k_1)A_k^2e^{-2\pi\abs k_1\sigma}.\]
  Then $\sum_kA_ke^{-2\pi\abs k_1\sigma}\le B^{\rm W}_\sigma$, and for $0<\sigma\le1$,
  \[B^{\rm W}_\sigma\le\frac{b_{\rm W}(\tau)}{\sigma^{2\tau+3}},\qquad
  b_{\rm W}(\tau)=2\bigl(C_0+C_1+C_{2\tau}+C_{2\tau+1}\bigr),\qquad C_0=\bigl(1+\tfrac1\pi\bigr)^2,\quad
  C_q=\frac4\pi\Bigl(\frac{q+1}{e\pi}\Bigr)^{q+1}\ (q>0),\]
  a sum of four terms counted with multiplicity (for $\tau=0$ the term $C_{2\tau}$ is $C_0$; for $\tau=\frac12$,
  $C_{2\tau}=C_1$ occurs twice). Since $n=2$, $2\tau+3=2\tau+n+1$.
\end{enumerate}
\end{lemma}

\begin{proof}
\proofstep{Counting}\label{stepW:w4:count}
For $N\ge1$, $k_1$ ranges over $-N,\dots,N$; $\abs{k_1}=N$ gives one point each and $\abs{k_1}<N$ gives two
($k_2=\pm(N-\abs{k_1})$), in total $2+2(2N-1)=4N$.
\proofstep{Product}\label{stepW:w4:coth}
$\sum_ke^{-2\pi\abs k_1\sigma}=\bigl(\sum_{m\in\Z}e^{-2\pi\abs m\sigma}\bigr)^2$ and
$\sum_{m\in\Z}e^{-x\abs m}=1+\frac{2}{e^x-1}=\coth\frac x2$ (\ref{bpW:w4:ext-exp}); $\coth t=1+\frac2{e^{2t}-1}\le1+\frac1t$.
\proofstep{Weighted sums}\label{stepW:w4:wsum}
For $q>0$ the term $k=0$ vanishes and, with $x=2\pi\sigma$, $S_q=4\sum_{N\ge1}N^{q+1}e^{-xN}$ by (i). By
\ref{bpW:w4:ext-exp}, $N^{q+1}e^{-xN/2}\le(2(q+1)/(ex))^{q+1}$ and $\sum_{N\ge1}e^{-xN/2}=1/(e^{x/2}-1)\le2/x$.
\proofstep{The constant}\label{stepW:w4:bw}
$A_k\ge1$ gives the first claim. The term $k=0$ of $B^{\rm W}_\sigma$ is $1$ ($A_0=1$). For $k\ne0$,
$(1+\abs k_1)A_k^2\le2(1+\abs k_1)(1+\abs k_1^{2\tau})=2\bigl(1+\abs k_1+\abs k_1^{2\tau}+\abs k_1^{2\tau+1}\bigr)$. Summing over
$k\ne0$, the four sums are at most $S_0-1$, $S_1$, $S_{2\tau}$, $S_{2\tau+1}$ (for $\tau=0$ the third is again at most
$S_0-1$), so $B^{\rm W}_\sigma\le1+2(S_0-1)+2(S_1+S'_{2\tau}+S_{2\tau+1})\le2(S_0+S_1+S'_{2\tau}+S_{2\tau+1})$, where
$S'_{2\tau}=S_{2\tau}$ for $\tau>0$ and $S'_0=S_0$. By (ii)--(iii), $S_0\le C_0\sigma^{-2}$ and $S_q\le C_q\sigma^{-(q+2)}$ for $q>0$,
each term counted separately even when two exponents coincide ($\tau=\frac12$), and
$\sigma^{-(q+2)}\le\sigma^{-(2\tau+3)}$ for $\sigma\le1$ and $q\le2\tau+1$.
\end{proof}
\source{ours. Checked numerically for $\tau\in\{0,0.5,1,3,3.5,5\}$, $\sigma\in[0.02,1]$
(\texttt{checks.py}); e.g.\ $b_{\rm W}(3)\approx5.76$.}

\subsection{Divisors and the multiplier solve}\label{bpW:w4:solve-sec}

\begin{lemma}[Divisor bounds]\label{bpW:w4:divisors}
Assume \eqref{bpW:w4:freq} and $(\mathrm H_\alpha)$. For every divisor $\delta$ of Definition~\ref{bpW:w4:melnikov}
with weight $w$, all $\xi,\zeta\in\Pi$ and $k\in\Z^2$:
\[\abs{\delta(\xi)}\ge\frac{\alpha w}{A_k},\qquad
\Bigl\lvert\frac1{\delta(\xi)}-\frac1{\delta(\zeta)}\Bigr\rvert\le\frac{2(1+\abs k_1)A_k^2M}{\alpha^2w}\abs{\xi-\zeta}.\]
The weights satisfy $w^Q_{ij}\ge\abs{i-j}$, $w^P_{ij}\ge\abs{i-j}$ for $i\ne j$, and all weights are $\ge1$.
\end{lemma}

\begin{proof}
\proofstep{Membership}\label{stepW:w4:member}
Each divisor is $\langle k,\omega\rangle+\langle l,\Omega\rangle$ with $l\in\{0,\mp e_j,\mp(e_i+e_j),e_i-e_j\}$, $\abs l\le2$,
and $(k,l)\ne0$ by the exclusions of Definition~\ref{bpW:w4:melnikov}; so $(k,l)\in\mathcal Z$. Then
$\jap l=1,\ j,\ i+j,\ \max(1,\abs{i-j})$ respectively, which is the weight, and $(\mathrm H_\alpha)$ gives the lower
bound.
\proofstep{Lipschitz}\label{stepW:w4:divlip}
$\abs{\delta(\xi)-\delta(\zeta)}\le\abs k_1M_\omega\abs{\xi-\zeta}+\sum_j\abs{l_j}\abs{\Omega_j(\xi)-\Omega_j(\zeta)}\le
(\abs k_1+2)M\abs{\xi-\zeta}$ by \eqref{bpW:w4:freq} ($\abs{\Omega_j(\xi)-\Omega_j(\zeta)}\le M_\Omega\abs{\xi-\zeta}/j$ and
$\abs l\le2$). Divide by $\abs{\delta(\xi)\delta(\zeta)}\ge\alpha^2w^2/A_k^2\ge\alpha^2w/A_k^2$ and use
$\abs k_1+2\le2(1+\abs k_1)$.
\end{proof}
\source{P96-K pp.~126--127. Deviations: (a) P96-K p.~127 asserts ``$\abs{\langle k,\omega\rangle+\Omega_i-\Omega_j}\ge
\alpha(1+\abs{i-j})/A_k$, since $d\ge1$''; from $\jap l_1=\max(1,\abs{i-j})$ only $\alpha\max(1,\abs{i-j})/A_k$ follows,
a factor up to $2$ smaller (suspected erratum; harmless, since Lemma~A.1 needs only $\abs{i-j}$). (b) For $d=1$ the
interpolation $2\abs{i^d-j^d}\ge\abs{i-j}(i^{d-1}+j^{d-1})$ and the factor $(i^\delta+j^\delta)$ of p.~127 are not needed:
$\abs{\Omega_i(\xi)-\Omega_i(\zeta)}\le M_\Omega/i$ is used directly.}

\begin{lemma}[Hadamard division of operators]\label{bpW:w4:schur}
Let $K_S=1+\pi/\sqrt3$ ($<2.82$). Let $d=(d_{ij})_{i,j\ge1}$ be complex numbers with $\abs{d_{ij}}\le c/\abs{i-j}$ for
$i\ne j$ and $\abs{d_{ii}}\le c$. Let $\mu,\nu$ be positive sequences and $A$ a matrix whose weighted matrix
$(\mu_iA_{ij}\nu_j)$ defines a bounded operator on $\ell^2$ of norm $N_A$. Then $(\mu_id_{ij}A_{ij}\nu_j)$ defines a
bounded operator of norm $\le cK_SN_A$. In particular:
\begin{enumerate}[label=(\roman*)]
\item ($\mu_i=\omega_{a,\bar p}(i+2)$, $\nu_j=\omega_{a,p}(j+2)^{-1}$.) The Hadamard map
  $\mathcal H_d:M\mapsto(d_{ij}M_{ij})$ is a bounded linear map $\mathcal L\to\mathcal L$ of norm $\le cK_S$, and
  $\mathcal H_d(M)^T=\mathcal H_{d^T}(M^T)$.
\item ($\mu_i=\nu_i=\omega_{a,p}(i+2)^{-1}$.) The same Hadamard map on bilinear forms, $\beta\mapsto
  \bigl((z,z')\mapsto\sum d_{ij}\beta(e_j,e_i)z_jz'_i\bigr)$, is bounded $\mathcal B\to\mathcal B$ with norm $\le cK_S$, and
  $\Lambda\circ\mathcal H_d=\mathcal H_d\circ\Lambda$.
\item For diagonal $d$ ($d_{ij}=0$ for $i\ne j$), $\abs{d_{ii}}\le c$, the bound improves to $c$, for every
  $\mu,\nu$; in particular $\abs{M_{ii}}\mu_i\nu_i\le\norm M$ for $M\in\mathcal L$, i.e.\ $(i+2)\abs{M_{ii}}\le\norm M$.
\end{enumerate}
\end{lemma}

\begin{proof}
\proofstep{Transport}\label{stepW:w4:schurtr}
Apply Lemma~\ref{bpW:w1:schur}(iii) to the $\ell^2$-matrix $\hat A=(\mu_iA_{ij}\nu_j)$ and $\hat F=(\mu_id_{ij}A_{ij}\nu_j)$:
$\abs{\hat F_{ij}}\le c\abs{\hat A_{ij}}/\abs{i-j}$ off the diagonal and $\abs{\hat F_{ii}}\le c\abs{\hat A_{ii}}$.
\proofstep{Representatives}\label{stepW:w4:schurrep}
In (i), the weighted matrix of $M\in\mathcal L$ is the matrix of its representative
$\tau_{a,\bar p}\circ M\circ\kappa_{a,p}$ (Definition~\ref{bpW:w1:transport}, restricted to the normal entries), so
bounded weighted matrices are exactly elements of $\mathcal L$ with the same norm; the transpose identity is
entrywise. In (ii), let $\hat\beta(x,y)=\beta(\kappa x,\kappa y)$ on $\ell^2_\C\times\ell^2_\C$, where $(\kappa x)_j=\mu_jx_j$ with
$\mu_j=\omega_{a,p}(j+2)^{-1}$ (the coefficient map of Definition~\ref{bpW:w1:transport} on normal entries); it is complex
bilinear with $\norm{\hat\beta}=\norm\beta_{\rm bil}$, since $\kappa$ is an isometry onto $\ell^N_{a,p}$. Let $C$ be
coordinatewise conjugation on $\ell^2_\C$ (a conjugate-linear isometry, $\langle Cu,w\rangle=\langle u,w\rangle_0$
by \ref{bpW:w1:ext-lp}). The form $B(v,w)=\hat\beta(Cv,w)$ is conjugate-linear in $v$, linear in $w$, bounded by
$\norm{\hat\beta}$; put $T=C\circ B^\sharp\circ C$ with $B^\sharp$ of \ref{bpW:w4:ext-riesz}. Then $T$ is complex linear,
continuous, and $\langle Tx,y\rangle_0=\langle CTx,y\rangle=\langle B^\sharp Cx,y\rangle=B(Cx,y)=\hat\beta(x,y)$. By the
norm formula of \ref{bpW:w4:ext-riesz} (applied to $CTx$, with $\norm{CTx}=\norm{Tx}$),
$\norm{Tx}=\sup_{\norm y\le1}\abs{\hat\beta(x,y)}$, so $\norm T=\norm{\hat\beta}=\norm\beta_{\rm bil}$. The matrix of $T$ is
$\langle Te_j,e_i\rangle_0=\hat\beta(e_j,e_i)=\mu_j\mu_i\beta(e_j,e_i)$, the weighted matrix. Conversely every bounded
operator $T$ on $\ell^2_\C$ defines the bounded form $(x,y)\mapsto\langle Tx,y\rangle_0$ of norm $\le\norm T$ by
\eqref{bpW:w4:pairing}-type Cauchy--Schwarz. So bounded weighted matrices correspond to elements of $\mathcal B$ with
the same norm, and Lemma~\ref{bpW:w1:schur}(iii) transfers; $\Lambda$ does not change entries.
\proofstep{Diagonal}\label{stepW:w4:schurdiag}
(iii) is Lemma~\ref{bpW:w1:schur}(ii) and Lemma~\ref{bpW:w1:multiplier}; $\mu_i\nu_i=(i+2)^{\bar p-p}=i+2$ in (i).
\end{proof}
\source{P96-K Lemma~A.1, pp.~145--146, and its use on p.~127 (``$\abs{\!\abs{\!\abs{\tilde F_k}\!}\!}\le3(A_k/\alpha)
\abs{\!\abs{\!\abs{\tilde R_k}\!}\!}$''); Lemma~\ref{bpW:w1:schur}. The bilinear version (ii) is ours; it is needed for the
$\theta$-derivatives of the quadratic components (Theorem~\ref{bpW:w4:lemma1}).}

\begin{lemma}[Multiplier solve]\label{bpW:w4:solve}
Let $E,E'$ be complex Banach spaces, $0<\sigma<s$, $\alpha>0$, $\kappa\ge1$, $M\ge0$, and for $k\in\Z^2$ and
$\xi\in\Pi$ let $T_k(\xi)\in L(E,E')$ satisfy
\[\norm{T_k(\xi)}\le\frac{\kappa A_k}\alpha,\qquad\norm{T_k(\xi)-T_k(\zeta)}\le\frac{2\kappa(1+\abs k_1)A_k^2M}{\alpha^2}\abs{\xi-\zeta}.\]
Let $f(\cdot,\xi):T_s\to E$ be complex differentiable, $1$-periodic, with $\norm{f(\theta,\xi)}\le m$ and
$\norm{f(\theta,\xi)-f(\theta,\zeta)}\le m^{\mathcal L}\abs{\xi-\zeta}$ on $T_s$. Define
\[(\mathcal Sf)(\theta,\xi)=\sum_{k\in\Z^2}T_k(\xi)\hat f_k(\xi)\,e^{2\pi i\langle k,\theta\rangle}.\]
\begin{enumerate}[label=(\roman*)]
\item The series converges absolutely and uniformly on $T_{s-\sigma}$; $\mathcal Sf(\cdot,\xi)$ is complex
  differentiable and $1$-periodic there, with coefficients $T_k\hat f_k$, and
  $\norm{\mathcal Sf(\theta,\xi)}\le\kappa B^{\rm W}_\sigma m/\alpha$.
\item $\norm{\mathcal Sf(\theta,\xi)-\mathcal Sf(\theta,\zeta)}\le\frac{\kappa B^{\rm W}_\sigma}{\alpha}\bigl(m^{\mathcal L}+\frac{2Mm}{\alpha}\bigr)\abs{\xi-\zeta}$ on $T_{s-\sigma}$.
\item If $\partial_{\theta_i}f$ is bounded on $T_s$, then $\partial_{\theta_i}\mathcal Sf=\mathcal S(\partial_{\theta_i}f)$
  on $T_{s-\sigma}$.
\item If $\Lambda:E'\to E''$ is continuous linear, then $\Lambda\mathcal Sf$ is the solve of $f$ with the operators
  $\Lambda T_k$. If $\Lambda_0:E\to E_0$, $\Lambda_1:E'\to E_1$ and $T'_k\in L(E_0,E_1)$ satisfy
  $\Lambda_1T_k=T'_k\Lambda_0$, then $\Lambda_1\mathcal Sf=\mathcal S'(\Lambda_0f)$.
\end{enumerate}
\end{lemma}

\begin{proof}
\proofstep{Coefficients}\label{stepW:w4:solvecoef}
By \ref{bpW:w4:F1}, $\norm{T_k\hat f_k}\le\frac{\kappa A_k}\alpha me^{-2\pi\abs k_1s}$, so
$\sum_k\norm{T_k\hat f_k}e^{2\pi\abs k_1(s-\sigma)}\le\frac{\kappa m}\alpha\sum_kA_ke^{-2\pi\abs k_1\sigma}\le\frac{\kappa m}\alpha B^{\rm W}_\sigma$
(Lemma~\ref{bpW:w4:lattice}(iv)); \ref{bpW:w4:F2} with $s'=s-\sigma$ gives (i).
\proofstep{Differences}\label{stepW:w4:solvediff}
$T_k(\xi)\hat f_k(\xi)-T_k(\zeta)\hat f_k(\zeta)=T_k(\xi)\bigl(\hat f_k(\xi)-\hat f_k(\zeta)\bigr)+\bigl(T_k(\xi)-T_k(\zeta)\bigr)\hat f_k(\zeta)$;
by Lemma~\ref{bpW:w4:fourier-ops}(iii) and \ref{bpW:w4:F1} its norm is at most
$\bigl(\frac{\kappa A_k}\alpha m^{\mathcal L}+\frac{2\kappa(1+\abs k_1)A_k^2Mm}{\alpha^2}\bigr)\abs{\xi-\zeta}e^{-2\pi\abs k_1s}$.
Multiply by $e^{2\pi\abs k_1(s-\sigma)}$, sum, and use $A_k\le(1+\abs k_1)A_k^2$ and the definition of $B^{\rm W}_\sigma$; the
difference of the two series is the series of differences, bounded as in \ref{bpW:w4:F2}.
\proofstep{Commutation}\label{stepW:w4:solvecomm}
(iii): by Lemma~\ref{bpW:w4:fourier-ops}(ii), $T_k\widehat{(\partial_{\theta_i}f)}_k=2\pi ik_iT_k\hat f_k$, and by
Step~\ref{stepW:w4:solvecoef} applied to the bounded map $\partial_{\theta_i}f$ these terms are summable with the
weight $e^{2\pi\abs k_1(s-\sigma)}$. So the termwise derivative of the series of $\mathcal Sf$ converges uniformly on
$T_{s-\sigma}$, and termwise differentiation (\ref{bpW:w1:ext-unifderiv}) gives
$\partial_{\theta_i}\mathcal Sf=\sum_k2\pi ik_iT_k\hat f_ke^{2\pi i\langle k,\theta\rangle}=\mathcal S(\partial_{\theta_i}f)$.
\proofstep{Linear images}\label{stepW:w4:solvelin}
(iv): $\Lambda$ commutes with absolutely convergent sums; and $\Lambda_1T_k\hat f_k=T'_k\Lambda_0\hat f_k=T'_k\widehat{(\Lambda_0f)}_k$
by Lemma~\ref{bpW:w4:fourier-ops}(i).
\end{proof}
\source{P96-K pp.~126--127 (``It follows that $\norm{\dot F}_{D(s-\sigma)}\le\sum_k\norm{\dot F_k}e^{\abs k(s-\sigma)}\le\dots$'',
``summing up as before''), in the Banach-valued form that replaces the square-sum step
(Remark~\ref{bpW:w4:counterexample}). Lean: the series is a \Lean{tsum} in $E'$, summable in norm by
Step~\ref{stepW:w4:solvecoef}.}

\subsection{The Lie derivative along the normal form}\label{bpW:w4:lie-sec}

At a fixed real parameter $\xi$ the normal form is formally $N=\langle\omega(\xi),y\rangle+\sum_j\Omega_j(\xi)z_j\bar z_j$.
Its field is unbounded on $\ell^N_{a,p}$ ($\Omega_j\sim j$). For general $a,p\ge0$ the series $\sum_j\Omega_jz_j\bar z_j$
need not even converge: for $a=p=0$, $\Omega_j=j+2$ and $z_j=\bar z_j=1/j$ it is $\sum_j(1/j+2/j^2)=\infty$. So we
define the bracket of $N$ with $F\in\mathcal T$ as a \emph{weak pairing} through the smoothing gradients of $F$, which
exists for all $a,p\ge0$, and identify it with an actual derivative of $N$ only under the guard
\begin{equation}\label{bpW:w4:guard}
\gamma_{a,p}:=\sup_{j\ge1}\frac{j+2}{\omega_{a,p}(j+2)^2}<\infty.\tag{G}
\end{equation}
(G) holds if and only if $a>0$ or $p\ge\frac12$: for $a>0$, $J^{1-2p}e^{-2aJ}\to0$; for $a=0$, $J^{1-2p}$ is bounded iff
$p\ge\frac12$. The wave consumer $a=p=1$ satisfies (G) with $\gamma_{1,1}=e^{-6}/3$ (at $J=3$).

\begin{lemma}[Lie derivative and weak bracket]\label{bpW:w4:lie}
Assume \eqref{bpW:w4:freq} and let $F\in\mathcal T_{s'}$ (Definition~\ref{bpW:w4:taylor}) with components
$a_F,b_F,c^\pm_F,P_F,Q_F,P'_F$. For $(\theta,y,z,\bar z)\in T_{s'}\times\C^2\times(\ell^N_{a,p})^2$ put
\[L_NF=\frac1{2\pi}\langle\omega,F_\theta\rangle-i\langle F_z,\Omega z\rangle_0+i\langle F_{\bar z},\Omega\bar z\rangle_0,\qquad
\langle\omega,F_\theta\rangle=\sum_i\omega_i\partial_{\theta_i}F,\]
and the \emph{weak bracket}
\[\langle\!\langle N,F\rangle\!\rangle=-\frac1{2\pi}\langle\omega,F_\theta\rangle+\langle\Omega\bar z,\dot Z\rangle_0+\langle\Omega z,\dot{\bar Z}\rangle_0,
\qquad(\dot Z,\dot{\bar Z})=(-iF_{\bar z},iF_z),\]
which is the formal expression for $DN\cdot X_F$: $-\frac1{2\pi}F_\theta$ is the $y$-component and $(\dot Z,\dot{\bar Z})$ the
normal components of $X_F$ by \eqref{bpW:w4:hameq}.
\begin{enumerate}[label=(\roman*)]
\item ($a,p\ge0$.) $\Omega z\in\ell^N_{a,p-1}$ with $\norm{\Omega z}_{a,p-1}\le(1+M_0)\norm z$; the pairings converge
  absolutely, $L_NF$ and $\langle\!\langle N,F\rangle\!\rangle$ are continuous in $(z,\bar z)$ and polynomials of degree $\le2$ in
  $(y,z,\bar z)$.
\item ($a,p\ge0$.) $\langle\!\langle N,F\rangle\!\rangle=-L_NF$.
\item (Coefficients, $a,p\ge0$.) For finitely supported $z,\bar z$, $L_NF$ is the finite expression obtained from
  \eqref{bpW:w4:R} for $F$ by replacing the coefficient of $e^{2\pi i\langle k,\theta\rangle}$ in
  $a_F$, $b_F$, $c^+_{F,j}$, $c^-_{F,j}$, $P_{F,ij}$, $Q_{F,ij}$, $P'_{F,ij}$ by $i\delta$ times it, with
  $\delta=\langle k,\omega\rangle$ (for $a_F$, $b_F$), $\delta^+_{k,j}$, $\delta^-_{k,j}$, $\delta^P_{k,ij}$, $\delta^Q_{k,ij}$
  (with $\delta^Q_{0,jj}=0$), $\delta^{P'}_{k,ij}$ respectively.
\item (Derivative, under (G).) If \eqref{bpW:w4:guard} holds, then $N_\Omega(z,\bar z)=\sum_j\Omega_jz_j\bar z_j$
  converges absolutely and is a continuous bilinear form on $\ell^N_{a,p}\times\ell^N_{a,p}$ with
  $\abs{N_\Omega(z,\bar z)}\le(1+M_0)\gamma_{a,p}\norm z\norm{\bar z}$; $N=\langle\omega,y\rangle+N_\Omega(z,\bar z)$ is a
  continuous polynomial on $\C^2\times\C^2\times(\ell^N_{a,p})^2$ (independent of $\theta$), Fr\'echet differentiable with
  $DN(w)(\Theta,Y,Z,\bar Z)=\langle\omega,Y\rangle+N_\Omega(Z,\bar z)+N_\Omega(z,\bar Z)$, and
  \[DN(w)\cdot X_F(w)=\langle\!\langle N,F\rangle\!\rangle(w)=-L_NF(w).\]
  Without (G) the left side is not defined in general (example above), and only (ii) is available.
\end{enumerate}
\end{lemma}

\begin{proof}
\proofstep{Bounds}\label{stepW:w4:liebd}
$\abs{\Omega_j}\le j+2+M_0/j\le(1+M_0)(j+2)$, so $\abs{\Omega_j}\omega_{a,p-1}(j+2)\le(1+M_0)\omega_{a,p}(j+2)$. The
gradients $F_z,F_{\bar z}$ lie in $\ell^N_{a,\bar p}=\ell^N_{a,p+1}$, and \eqref{bpW:w4:pairing} with $q=p-1$, $q'=p+1$
($q+q'=2p\ge0$) bounds the pairings. Continuity in $(z,\bar z)$: $F_z$, $F_{\bar z}$ are affine continuous in
$(z,\bar z)$ and the pairing is a bounded bilinear form on $\ell^N_{a,p-1}\times\ell^N_{a,p+1}$.
\proofstep{Weak bracket}\label{stepW:w4:liebr}
Substitute $(\dot Z,\dot{\bar Z})$: $\langle\!\langle N,F\rangle\!\rangle=-\frac1{2\pi}\langle\omega,F_\theta\rangle-i\langle\Omega\bar z,F_{\bar z}\rangle_0
+i\langle\Omega z,F_z\rangle_0$, which is $-L_NF$ by symmetry of $\langle\cdot,\cdot\rangle_0$. This is an identity of explicit
expressions; no derivative of $N$ is involved.
\proofstep{Coefficients}\label{stepW:w4:liecoef}
For finitely supported $z,\bar z$ every pairing is a finite sum. Write $F_z=c^+_F+2P_Fz+Q_F^T\bar z$,
$F_{\bar z}=c^-_F+Q_Fz+2P'_F\bar z$ (Definition~\ref{bpW:w4:taylor}) and collect: the coefficient of $z_j$ is
$\frac1{2\pi}\langle\omega,\partial_\theta c^+_{F,j}\rangle-i\Omega_jc^+_{F,j}$; of $z_iz_j$ (symmetrized),
$\frac1{2\pi}\langle\omega,\partial_\theta P_{F,ij}\rangle-i(\Omega_i+\Omega_j)P_{F,ij}$; of $z_j\bar z_i$,
$\frac1{2\pi}\langle\omega,\partial_\theta Q_{F,ij}\rangle-i(\Omega_j-\Omega_i)Q_{F,ij}$; and similarly for $c^-_F$, $P'_F$ with
$+i\Omega$. By Lemma~\ref{bpW:w4:fourier-ops}, the $k$-th coefficient of $\frac1{2\pi}\langle\omega,\partial_\theta g\rangle$
is $i\langle k,\omega\rangle\hat g_k$, and the matrix entries are continuous linear functionals of the operators.
\proofstep{The guarded derivative}\label{stepW:w4:lieder}
Under (G), $\abs{\Omega_jz_j\bar z_j}\le(1+M_0)(j+2)\abs{z_j}\abs{\bar z_j}\le(1+M_0)\gamma_{a,p}\,\omega_{a,p}(j+2)^2\abs{z_j}\abs{\bar z_j}$,
and Cauchy--Schwarz gives absolute convergence and the bound; so $N_\Omega$ is a bounded bilinear form and $N$ a
continuous polynomial, whose derivative is the displayed one (Mathlib: \Lean{ContinuousLinearMap.hasFDerivAt} for the
linear part and the bilinear-map derivative \Lean{IsBoundedBilinearMap.hasFDerivAt}). The direction $X_F(w)$ has
normal components $\dot Z,\dot{\bar Z}\in\ell^N_{a,\bar p}\subset\ell^N_{a,p}$ (Lemma~\ref{bpW:w1:embed}) and $y$-component
$-\frac1{2\pi}F_\theta$, so $DN(w)\cdot X_F(w)=-\frac1{2\pi}\langle\omega,F_\theta\rangle+N_\Omega(\dot Z,\bar z)+N_\Omega(z,\dot{\bar Z})$,
and $N_\Omega(\dot Z,\bar z)=\langle\Omega\bar z,\dot Z\rangle_0$, $N_\Omega(z,\dot{\bar Z})=\langle\Omega z,\dot{\bar Z}\rangle_0$ termwise
(both series converge absolutely). This is $\langle\!\langle N,F\rangle\!\rangle$, and (ii) concludes.
\end{proof}
\source{P96-K p.~125 (``by comparison of coefficients \dots\ $iF_{kmq\bar q}=R_{kmq\bar q}/(\langle k,\omega\rangle+\langle q-\bar q,\Omega\rangle)$''),
in the W0 sign convention: our $\langle\cdot,\Omega\rangle$-terms are those of P96-K with $q\leftrightarrow\bar q$
(Remark~\ref{bpW:w4:melnikov-check}(b)). P96-K writes the bracket as $\{F,N\}$; here $\{F,N\}:=L_NF=-\langle\!\langle N,F\rangle\!\rangle$,
which equals $-DN\cdot X_F$ under (G). P96-K takes $a\ge0$, $p\ge0$ (p.~120) and does not address the convergence of
$N$; the guard (G) is ours. The identity (iv) is the input W5 needs to expand $N\circ X^t_F$; W5 must use it in this
form because $X_N$ is unbounded.}

\subsection{Reading the components off the field norm}\label{bpW:w4:extract-sec}

\begin{lemma}[Coefficient extraction]\label{bpW:w4:extract}
Let $E$ be a complex Banach space, $\rho_1,\rho_2>0$, and $g(\zeta,\eta)=\sum_{p+q\le2}g_{pq}\zeta^p\eta^q$ with
$g_{pq}\in E$. If $\norm{g(\zeta,\eta)}\le K$ whenever $\abs\zeta<\rho_1$, $\abs\eta<\rho_2$, then
$\norm{g_{pq}}\le K\rho_1^{-p}\rho_2^{-q}$ for all $p+q\le2$. The same holds for one variable.
\end{lemma}

\begin{proof}
For $0<\rho'_i<\rho_i$,
$g_{pq}\rho_1'^{\,p}\rho_2'^{\,q}=\frac1{16}\sum_{u,v=0}^3i^{-up-vq}g(\rho'_1i^u,\rho'_2i^v)$: expanding $g$, the
monomial $(p',q')$ contributes $\frac1{16}\sum_{u,v}i^{u(p'-p)+v(q'-q)}g_{p'q'}\rho_1'^{\,p'}\rho_2'^{\,q'}$, and since
$\abs{p'-p},\abs{q'-q}\le2<4$ this vanishes unless $(p',q')=(p,q)$ (\ref{bpW:w4:ext-roots}). The triangle inequality
gives $\norm{g_{pq}}\le K\rho_1'^{-p}\rho_2'^{-q}$; let $\rho'_i\uparrow\rho_i$.
\end{proof}
\source{ours; it replaces the generalized Cauchy estimate of P96-K Lemma~A.3 (p.~147) for the polynomial
dependence on $(y,z,\bar z)$.}

\begin{lemma}[Component bounds]\label{bpW:w4:components}
Let $s,r>0$, $R\in\mathcal T_s$, and suppose $\abs{X_R(\theta,y,z,\bar z)}_r\le\varepsilon$ on $D(s,r)$. Then for all
$\theta\in T_s$ and $i\in\{1,2\}$:
\begin{enumerate}[label=(\roman*)]
\item $\abs{b(\theta)}\le2\pi\varepsilon$, $\norm{c^\pm(\theta)}_{\bar p}\le r\varepsilon/\sqrt2$,
  $\norm{2P(\theta)},\norm{Q(\theta)},\norm{Q(\theta)^T},\norm{2P'(\theta)}\le\varepsilon$;
\item $\abs{\partial_{\theta_i}a}\le2\pi r^2\varepsilon$, $\abs{\partial_{\theta_i}b}_1\le2\pi\varepsilon$,
  $\norm{\partial_{\theta_i}c^\pm}_*\le2\sqrt2\pi r\varepsilon$, $\norm{\Lambda\partial_{\theta_i}P}_{\rm bil},
  \norm{\Lambda\partial_{\theta_i}P'}_{\rm bil}\le4\pi\varepsilon$, $\norm{\Lambda\partial_{\theta_i}Q}_{\rm bil}\le16\pi\varepsilon$.
\end{enumerate}
\end{lemma}

\begin{proof}
\proofstep{Test points}\label{stepW:w4:testpts}
By \eqref{bpW:w4:zu}, the points $(\theta,y,0,0)$ with $\abs y<r^2$; $(\theta,0,z,0)$ and $(\theta,0,0,\bar z)$ with
$\norm z,\norm{\bar z}<r/\sqrt2$; and $(\theta,0,z,\bar z)$ with $\norm z,\norm{\bar z}<r/(2\sqrt2)$ lie in $D(s,r)$
(for the second, $u=z/\sqrt2$, $v=z/(i\sqrt2)$ and $\norm u+\norm v=\sqrt2\norm z$). On $D(s,r)$ the hypothesis gives
$\abs{R_y}\le2\pi\varepsilon$, $\abs{R_\theta}\le2\pi r^2\varepsilon$, and
$\norm{R_z}_{\bar p},\norm{R_{\bar z}}_{\bar p}\le(\norm{R_u}_{\bar p}+\norm{R_v}_{\bar p})/\sqrt2\le r\varepsilon/\sqrt2$.
\proofstep{Zeroth and first order}\label{stepW:w4:first}
$R_y=b$ and $R_z(\theta,0,0,0)=c^+$, $R_{\bar z}(\theta,0,0,0)=c^-$. For $\norm z=1$,
$\zeta\mapsto R_z(\theta,0,\zeta z,0)=c^++2\zeta Pz$ is bounded by $r\varepsilon/\sqrt2$ on $\abs\zeta<r/\sqrt2$, so
$\norm{2Pz}_{\bar p}\le\varepsilon$ by Lemma~\ref{bpW:w4:extract}. Likewise $R_z(\theta,0,0,\eta\bar z)=c^++\eta Q^T\bar z$,
$R_{\bar z}(\theta,0,\zeta z,0)=c^-+\zeta Qz$, $R_{\bar z}(\theta,0,0,\eta\bar z)=c^-+2\eta P'\bar z$.
\proofstep{Angle derivatives}\label{stepW:w4:angle}
$\partial_{\theta_i}R(\theta,y,0,0)=\partial_{\theta_i}a+\langle\partial_{\theta_i}b,y\rangle$: its value at $y=0$ and, for
$\abs y=1$, the coefficient of $\zeta$ in $\zeta\mapsto\partial_{\theta_i}R(\theta,\zeta y,0,0)$ ($\abs\zeta<r^2$) give the
first two bounds ($\abs\cdot_1$ is the dual of the max norm). For $\norm z=1$,
$\zeta\mapsto\partial_{\theta_i}R(\theta,0,\zeta z,0)=\partial_{\theta_i}a+\zeta\langle\partial_{\theta_i}c^+,z\rangle_0+\zeta^2\langle\partial_{\theta_i}Pz,z\rangle_0$
on $\abs\zeta<r/\sqrt2$ gives $\abs{\langle\partial_{\theta_i}c^+,z\rangle_0}\le2\pi r^2\varepsilon\sqrt2/r$ and
$\abs{\langle\partial_{\theta_i}Pz,z\rangle_0}\le2\pi r^2\varepsilon\cdot2/r^2$. For $\norm z=\norm{\bar z}=1$, the coefficient of
$\zeta\eta$ in $\partial_{\theta_i}R(\theta,0,\zeta z,\eta\bar z)$ on $\abs\zeta,\abs\eta<r/(2\sqrt2)$ is
$\langle\partial_{\theta_i}Qz,\bar z\rangle_0$, bounded by $2\pi r^2\varepsilon\cdot8/r^2$. The $c^-$, $P'$ cases are symmetric.
\proofstep{Polarization}\label{stepW:w4:polar}
For symmetric $\beta=\Lambda\partial_{\theta_i}P$ and unit $z,z'$:
$\beta(z,z')=\frac14\bigl(\beta(z+z',z+z')-\beta(z-z',z-z')\bigr)$, so
$\abs{\beta(z,z')}\le\frac14\cdot4\pi\varepsilon\,(\norm{z+z'}^2+\norm{z-z'}^2)=4\pi\varepsilon$ by
\ref{bpW:w4:ext-parallelogram}. ($\partial_{\theta_i}P$ is symmetric because $P$ is.)
\end{proof}
\source{P96-K p.~126 (``$\norm{\dot R}_{D(s)}\le r\abs{X_R}_r$''), p.~127 (``$\abs{\!\abs{\!\abs{\tilde R}\!}\!}_{D(s)}\le\frac1r\norm{R_z}_{D(s,r)}\le\abs{X_R}_r$'')
and p.~128 (``All the other components $F^{ij}$ admit the same estimates''); the bounds for the angle derivatives are
not in the source, which leaves the $y$-component of $X_F$ implicit.}

\subsection{The homological equation}\label{bpW:w4:main-sec}

\begin{theorem}[P96-K Lemma~1 for the wave setting]\label{bpW:w4:lemma1}
Assume Convention~\ref{bpW:w4:index}, \eqref{bpW:w4:freq} on $\Pi\subset\R^2$ with $M=M_\omega+M_\Omega$, $\tau\ge0$, and
$(\mathrm H_\alpha)$ for some $\alpha>0$. Let $0<\sigma<s$, $r>0$, $\varepsilon,\varepsilon^{\mathcal L}\ge0$, and for each
$\xi\in\Pi$ let $R(\cdot,\xi)\in\mathcal T_s$ satisfy, at every point of $D(s,r)$ and for all $\xi,\zeta\in\Pi$,
\[\abs{X_{R(\xi)}}_r\le\varepsilon,\qquad\abs{X_{R(\xi)}-X_{R(\zeta)}}_r\le\varepsilon^{\mathcal L}\abs{\xi-\zeta}.\]
Define $F(\cdot,\xi)$ by the components (all series in $\theta$, $e_k=e^{2\pi i\langle k,\theta\rangle}$)
\[a_F=\sum_{k\ne0}\frac{\hat a_k}{i\delta^0_k}e_k,\quad b_F=\sum_{k\ne0}\frac{\hat b_k}{i\delta^0_k}e_k,\quad
(c^\pm_F)_j=\sum_k\frac{(\hat c^\pm_k)_j}{i\delta^\pm_{k,j}}e_k,\]
\[(P_F)_{ij}=\sum_k\frac{(\hat P_k)_{ij}}{i\delta^P_{k,ij}}e_k,\quad(Q_F)_{ij}=\sum_{(k,i)\ne(0,j)}\frac{(\hat Q_k)_{ij}}{i\delta^Q_{k,ij}}e_k,\quad
(P'_F)_{ij}=\sum_k\frac{(\hat P'_k)_{ij}}{i\delta^{P'}_{k,ij}}e_k,\]
(the vector and operator series converging in $\ell^N_{a,\bar p}$ and $\mathcal L$, entries as displayed), and
\[\hat N=\langle\hat\omega,y\rangle+\sum_j\hat\Omega_jz_j\bar z_j,\qquad\hat\omega=[b],\quad\hat\Omega_j=([Q])_{jj}.\]
Let $C_1=7+10K_S$ ($<36$) and $B=B^{\rm W}_\sigma$ (Lemma~\ref{bpW:w4:lattice}). Then for every $\xi\in\Pi$:
\begin{enumerate}[label=(\roman*)]
\item $F(\cdot,\xi)\in\mathcal T_{s-\sigma}$, with $\hat a_{F,0}=0$, $\hat b_{F,0}=0$, $(\hat Q_{F,0})_{jj}=0$; if $R$ is real,
  so are $F$ and $\hat N$ ($\hat\omega\in\R^2$, $\hat\Omega_j\in\R$).
\item $L_NF+\hat N=R-\hat a_0$ on $T_{s-\sigma}\times\C^2\times(\ell^N_{a,p})^2$, with $L_N$ of Lemma~\ref{bpW:w4:lie};
  equivalently $R+\langle\!\langle N,F\rangle\!\rangle=\hat N+\hat a_0$ (weak bracket of Lemma~\ref{bpW:w4:lie}), and, under
  the guard \eqref{bpW:w4:guard} (in particular for $a=p=1$), $R+DN\cdot X_F=\hat N+\hat a_0$ with the Fr\'echet derivative of $N$.
\item On $D(s-\sigma,r)$: $\abs{X_F}_r\le\frac{C_1B}{\alpha}\varepsilon$ and
  $\abs{X_{F(\xi)}-X_{F(\zeta)}}_r\le\frac{C_1B}{\alpha}\bigl(\varepsilon^{\mathcal L}+\frac{2M}{\alpha}\varepsilon\bigr)\abs{\xi-\zeta}$.
\item $\abs{\hat\omega}\le2\pi\varepsilon$, $\sup_j(j+2)\abs{\hat\Omega_j}\le\varepsilon$, and the same with
  $\varepsilon^{\mathcal L}\abs{\xi-\zeta}$ for the differences; on $D(s,r)$, $\abs{X_{\hat N}}_r\le2\varepsilon$ and
  $\abs{X_{\hat N(\xi)}-X_{\hat N(\zeta)}}_r\le2\varepsilon^{\mathcal L}\abs{\xi-\zeta}$.
\end{enumerate}
\end{theorem}

\begin{proof}
\proofstep{Division operators}\label{stepW:w4:divops}
For each component define $T_k(\xi)$: division by $i\delta^0_k$ on $\C$ and $\C^2$ (zero for $k=0$); the diagonal
multipliers $x\mapsto(x_j/(i\delta^\pm_{k,j}))_j$ on $\ell^N_{a,\bar p}$ and on $\ell^N_*$; the Hadamard maps
$\mathcal H_d$ with $d_{ij}=1/(i\delta_{k,ij})$ for $\delta\in\{\delta^P,\delta^Q,\delta^{P'}\}$ ($d_{jj}=0$ when $k=0$ for
$\delta^Q$) on $\mathcal L$ and on $\mathcal B$. By Lemma~\ref{bpW:w4:divisors}, $\abs{d_{ij}}\le A_k/(\alpha w_{ij})$ with
$w_{ij}\ge\max(1,\abs{i-j})$, and the differences $d(\xi)-d(\zeta)$ satisfy the same with
$2(1+\abs k_1)A_k^2M\abs{\xi-\zeta}/\alpha^2$ in place of $A_k/\alpha$. Lemma~\ref{bpW:w4:schur} (and
Lemma~\ref{bpW:w1:multiplier} for the diagonal cases) gives the hypotheses of Lemma~\ref{bpW:w4:solve} with $\kappa=1$
for $a,b,c^\pm$ and $\kappa=K_S$ for $P,Q,Q^T,P'$.
\proofstep{The solution is in the class}\label{stepW:w4:class}
Each component of $F$ is $\mathcal Sg$ for the corresponding component $g$ of $R$ (Lemma~\ref{bpW:w4:solve}(i)); for
$Q_F^T$ use $\mathcal H_d(M)^T=\mathcal H_{d^T}(M^T)$ (Lemma~\ref{bpW:w4:schur}(i)), where $d^T$ is the division by
$\delta^Q_{k,ji}$, which obeys the same bounds. $P_F$, $P'_F$ are symmetric since $\delta^P_{k,ij}$, $\delta^{P'}_{k,ij}$ are
symmetric in $(i,j)$. The vanishing coefficients are the excluded terms.
\proofstep{The equation}\label{stepW:w4:equation}
Fix $y$ and finitely supported $z,\bar z$. Then $\theta\mapsto L_NF+\hat N-R+\hat a_0$ is complex differentiable,
periodic and bounded on $T_{s-\sigma}$ (the $\theta$-derivatives of the components of $F$ are bounded there by
Lemma~\ref{bpW:w4:solve}(iii) and Lemma~\ref{bpW:w4:components}(ii)). By Lemma~\ref{bpW:w4:lie}(iii) and the
definition of $F$, its $k$-th coefficient is a finite sum of terms $i\delta\cdot\hat g_k/(i\delta)-\hat g_k=0$ for the
nonresonant terms, while the resonant terms are $-\hat a_0+\hat a_0$, $\langle\hat\omega-\hat b_0,y\rangle$ and
$(\hat\Omega_j-(\hat Q_0)_{jj})z_j\bar z_j$, all zero. By \ref{bpW:w4:F3} the function vanishes. Both sides are
continuous in $(z,\bar z)\in(\ell^N_{a,p})^2$ (Lemma~\ref{bpW:w4:lie}(i); $\hat N$ by (iv) below), and finitely
supported sequences are dense (Lemma~\ref{bpW:w1:space}(iii)). The second form is Lemma~\ref{bpW:w4:lie}(ii), the third
Lemma~\ref{bpW:w4:lie}(iv).
\proofstep{Component bounds of $F$}\label{stepW:w4:Fcomp}
By Lemma~\ref{bpW:w4:components}, Step~\ref{stepW:w4:divops} and Lemma~\ref{bpW:w4:solve}(i),(iii),(iv), on $T_{s-\sigma}$:
$\abs{b_F}\le2\pi\varepsilon B/\alpha$; $\norm{c^\pm_F}_{\bar p}\le\frac{B}{\alpha}\frac{r\varepsilon}{\sqrt2}$;
$\norm{2P_F},\norm{Q_F},\norm{Q^T_F},\norm{2P'_F}\le K_S\frac{B}{\alpha}\varepsilon$; and for $i=1,2$,
$\abs{\partial_{\theta_i}a_F}\le\frac B\alpha2\pi r^2\varepsilon$, $\abs{\partial_{\theta_i}b_F}_1\le\frac B\alpha2\pi\varepsilon$,
$\norm{\partial_{\theta_i}c^\pm_F}_*\le\frac B\alpha2\sqrt2\pi r\varepsilon$,
$\norm{\Lambda\partial_{\theta_i}P_F}_{\rm bil},\norm{\Lambda\partial_{\theta_i}P'_F}_{\rm bil}\le K_S\frac B\alpha4\pi\varepsilon$,
$\norm{\Lambda\partial_{\theta_i}Q_F}_{\rm bil}\le K_S\frac B\alpha16\pi\varepsilon$. (For the bilinear bounds first apply
Lemma~\ref{bpW:w4:solve}(iv) with $\Lambda_0=\Lambda_1=\Lambda$ and the Hadamard maps on $\mathcal B$
(Lemma~\ref{bpW:w4:schur}(ii)), so $\Lambda P_F$ is the $\mathcal B$-valued solve of $\Lambda P$, then
Lemma~\ref{bpW:w4:solve}(iii) in $\mathcal B$, where $\partial_{\theta_i}\Lambda P=\Lambda\partial_{\theta_i}P$ is bounded by
Lemma~\ref{bpW:w4:components}(ii). For $c^\pm$ likewise with the inclusion $\ell^N_{a,\bar p}\subset\ell^N_*$,
$\norm c_*\le\norm c_{\bar p}$ by \eqref{bpW:w4:pairing}. For $\abs\cdot_1$ on $\C^2$, division by a scalar $i\delta^0_k$
multiplies the norm by at most $A_k/\alpha$.)
\proofstep{Assembly}\label{stepW:w4:assembly}
Let $(\theta,y,z,\bar z)\in D(s-\sigma,r)$, so $\abs y<r^2$ and $\norm z,\norm{\bar z}<r/\sqrt2$ by \eqref{bpW:w4:zu}.
\emph{$\theta$-component:} $\frac1{2\pi}\abs{F_y}=\frac1{2\pi}\abs{b_F}\le\frac B\alpha\varepsilon$.
\emph{$y$-component:} $\partial_{\theta_i}F=\partial_{\theta_i}a_F+\langle\partial_{\theta_i}b_F,y\rangle+\sum_\pm\langle\partial_{\theta_i}c^\pm_F,\cdot\rangle_0
+\Lambda\partial_{\theta_i}P_F(z,z)+\Lambda\partial_{\theta_i}Q_F(z,\bar z)+\Lambda\partial_{\theta_i}P'_F(\bar z,\bar z)$, so
\[\abs{\partial_{\theta_i}F}\le\frac B\alpha\varepsilon\,2\pi r^2\Bigl(1+1+2\cdot1+K_S\bigl(1+1+4\bigr)\Bigr)
=\frac B\alpha\varepsilon\,2\pi r^2(4+6K_S),\]
using $\abs{\langle\beta,y\rangle}\le\abs\beta_1\abs y$, $\norm z^2<r^2/2$ and $\norm z\norm{\bar z}<r^2/2$; hence
$\frac1{2\pi r^2}\abs{F_\theta}\le(4+6K_S)\frac B\alpha\varepsilon$.
\emph{Normal components:} $F_z=c^+_F+2P_Fz+Q_F^T\bar z$ gives
$\norm{F_z}_{\bar p}\le\frac B\alpha\frac{r\varepsilon}{\sqrt2}(1+2K_S)$, and the same for $F_{\bar z}$; by the formulas for
$F_u,F_v$, $\frac1r(\norm{F_u}_{\bar p}+\norm{F_v}_{\bar p})\le\frac{\sqrt2}r(\norm{F_z}_{\bar p}+\norm{F_{\bar z}}_{\bar p})\le
2(1+2K_S)\frac B\alpha\varepsilon$. Summing, $\abs{X_F}_r\le(1+4+6K_S+2+4K_S)\frac B\alpha\varepsilon=C_1\frac B\alpha\varepsilon$.
\proofstep{Lipschitz}\label{stepW:w4:lip}
$R(\xi)-R(\zeta)\in\mathcal T_s$ and $X_{R(\xi)-R(\zeta)}=X_{R(\xi)}-X_{R(\zeta)}$, so Lemma~\ref{bpW:w4:components} applied to
it gives every bound of Step~\ref{stepW:w4:Fcomp}'s inputs with $\varepsilon^{\mathcal L}\abs{\xi-\zeta}$ in place of
$\varepsilon$. Lemma~\ref{bpW:w4:solve}(ii) then gives each component difference of $F$ with the factor
$\varepsilon^{\mathcal L}+2M\varepsilon/\alpha$ in place of $\varepsilon$, and the assembly of Step~\ref{stepW:w4:assembly},
which is linear in the components, gives (iii).
\proofstep{The normal form correction}\label{stepW:w4:Nhat}
$\abs{\hat\omega}=\abs{[b]}\le\sup_{\R^2}\abs b\le2\pi\varepsilon$. $\hat\Omega_j=([Q])_{jj}$ and $\norm{[Q]}\le\sup\norm Q\le\varepsilon$
(the mean is an average), so $(j+2)\abs{\hat\Omega_j}\le\varepsilon$ by Lemma~\ref{bpW:w4:schur}(iii). For the field: since
$z_j\bar z_j=\frac12(u_j^2+v_j^2)$, $\hat N_u=\hat\Omega u$, $\hat N_v=\hat\Omega v$ and
$\norm{\hat\Omega u}_{\bar p}\le\sup_j(j+2)\abs{\hat\Omega_j}\norm u$, so $X_{\hat N}=(\frac1{2\pi}\hat\omega,0,\hat\Omega v,-\hat\Omega u)$ has
$\abs{X_{\hat N}}_r\le\varepsilon+\varepsilon(\norm u+\norm v)/r\le2\varepsilon$ on $D(s,r)$. Differences: the same with
$R(\xi)-R(\zeta)$.
\proofstep{Reality}\label{stepW:w4:real}
For real $R$ and real $\xi$, the reality relations of Definition~\ref{bpW:w4:taylor} read
$\hat a_{-k}=\overline{\hat a_k}$, $\hat b_{-k}=\overline{\hat b_k}$, $\hat c^-_{k}=\overline{\hat c^+_{-k}}$,
$\hat P'_k=\overline{\hat P_{-k}}$, $(\hat Q_k)_{ji}=\overline{(\hat Q_{-k})_{ij}}$. The divisors are real and satisfy
$\delta^0_{-k}=-\delta^0_k$, $\delta^-_{k,j}=-\delta^+_{-k,j}$, $\delta^{P'}_{k,ij}=-\delta^P_{-k,ij}$,
$\delta^Q_{k,ji}=-\delta^Q_{-k,ij}$; since $\overline{1/(i\delta)}=-1/(i\delta)$ for real $\delta$, the coefficients of $F$
satisfy the same relations. $\hat\omega$ and $\hat\Omega_j$ are means of functions real on $\R^2$ ($Q_{jj}$ is real there since
$Q^T=\overline Q$).
\end{proof}
\source{P96-K Lemma~1, p.~125, and its proof pp.~125--128. Deviations, all reported: (a) the operator-valued
summation of p.~127 is replaced (Remark~\ref{bpW:w4:counterexample}); (b) the constant $B_\sigma$ with
$B_\sigma^2=2^n\sum_k(1+\abs k)^2A_k^4e^{-2\abs k\sigma}$ is replaced by $B^{\rm W}_\sigma$; (c) constants $C_1<36$ and
$2C_1$ instead of $2^6$, $2^8$; (d) $\abs{X_{\hat N}}_r\le2\abs{X_R}_r$ instead of $\le\abs{X_R}_r$; (e) the equation holds
modulo the constant $\hat a_0$ (P96-K's $\hat N=[R]$, summed over $\abs m+\abs q=1$, omits the constant term of $R$, so
$\{F,N\}+\hat N=R$ holds only up to that constant; harmless, constants do not enter $X$); (f) $\theta$ has period one
and the $\theta$-component of $X$ carries the factor $(2\pi)^{-1}$.}
\sourcegap{P96-K p.~128 states that the components other than $F^{10}$, $F^{11}$ ``admit the same estimates, or even
better ones'' and that ``the estimates of $X_F$ thus follow'', without the bookkeeping of the $y$-component of $X_F$
(which contains $\partial_\theta F^{ij}$ for all components) or of the constants; Steps~\ref{stepW:w4:Fcomp}
and~\ref{stepW:w4:assembly} supply it.}

\begin{corollary}[P96-K Lemma~2]\label{bpW:w4:lemma2}
In the setting of Theorem~\ref{bpW:w4:lemma1} with $M>0$, let $0\le\lambda\le\alpha/M$ and write
$\abs{X}^\lambda=\varepsilon_X+\lambda\varepsilon_X^{\mathcal L}$ for a field with uniform bound $\varepsilon_X$ and Lipschitz
constant $\varepsilon_X^{\mathcal L}$ on the relevant domain. Then
\[\abs{X_{\hat N}}^\lambda_{r,D(s,r)}\le2\abs{X_R}^\lambda_{r,D(s,r)},\qquad
\abs{X_F}^\lambda_{r,D(s-\sigma,r)}\le\frac{3C_1B^{\rm W}_\sigma}{\alpha}\abs{X_R}^\lambda_{r,D(s,r)},\]
and for $0<\sigma\le1$, $B^{\rm W}_\sigma\le b_{\rm W}(\tau)\,\sigma^{-(2\tau+n+1)}$ ($n=2$), the bound (4) of P96-K with
$b=b_{\rm W}(\tau)$; the constant $3C_1<108$ plays the role of P96-K's absolute constant $a$.
\end{corollary}

\begin{proof}
By Theorem~\ref{bpW:w4:lemma1}(iii),
$\abs{X_F}^\lambda\le\frac{C_1B}\alpha\bigl(\varepsilon+\lambda\varepsilon^{\mathcal L}+\frac{2\lambda M}\alpha\varepsilon\bigr)\le\frac{3C_1B}\alpha(\varepsilon+\lambda\varepsilon^{\mathcal L})$
since $\lambda M/\alpha\le1$; (iv) gives the first bound; Lemma~\ref{bpW:w4:lattice}(iv) the last.
\end{proof}
\source{P96-K Lemma~2 and (4), p.~128.}

\subsection{Comparison with P96-K and consequences for W5--W6}\label{bpW:w4:compare-sec}

\begin{remark}[Constants, strip powers, smallness exponents]\label{bpW:w4:compare}
\begin{center}\small
\begin{tabular}{@{}lll@{}}
\toprule
Quantity & P96-K (pp.~125, 128) & Here (period-one units, $n=2$)\\
\midrule
Coefficient summation & $L^2$ bound and Cauchy--Schwarz & Banach bound \ref{bpW:w4:F1} and $\ell^1$ sum\\
Strip constant & $B_\sigma^2=2^n\sum_k(1+\abs k)^2A_k^4e^{-2\abs k\sigma_K}$ & $B^{\rm W}_\sigma=\sum_k(1+\abs k_1)A_k^2e^{-2\pi\abs k_1\sigma}$\\
Power of $\sigma$ & $B_\sigma\le b\,\sigma^{-(2\tau+n+1)}$ & $B^{\rm W}_\sigma\le b_{\rm W}(\tau)\sigma^{-(2\tau+n+1)}$ ($\sigma\le1$)\\
$\abs{X_F}_{r,D(s-\sigma,r)}$ & $2^6B_\sigma\abs{X_R}_r/\alpha$ & $C_1B^{\rm W}_\sigma\varepsilon/\alpha$, $C_1<36$\\
$\abs{X_F}^{\mathcal L}$ & $2^8\frac{B_\sigma}\alpha(\abs{X_R}^{\mathcal L}+\frac M\alpha\abs{X_R})$ & $C_1\frac{B^{\rm W}_\sigma}\alpha(\varepsilon^{\mathcal L}+\frac{2M}\alpha\varepsilon)$\\
$\abs{X_{\hat N}}$ & $\le\abs{X_R}$ & $\le2\varepsilon$\\
Lemma~2 constant $a$ & absolute & $3C_1<108$\\
\bottomrule
\end{tabular}
\end{center}
(i) The Banach route costs $\sigma^{-n}$ in the lattice sum where the $L^2$ route costs $\sigma^{-n/2}$; but P96-K
gains nothing from this, since it bounds $B_\sigma$ by $\sigma^{-(2\tau+n+1)}$ while its own definition gives
$\sigma^{-(2\tau+1+n/2)}$. Our sum needs only the weight $(1+\abs k_1)A_k^2$ (not $(1+\abs k)^2A_k^4$) because the
$\theta$-derivatives of $F$ are solved from the $\theta$-derivatives of $R$ (Step~\ref{stepW:w4:Fcomp}), so it attains the
same power $\sigma^{-(2\tau+n+1)}$.
(ii) Hence (8) of P96-K p.~129 holds with the same exponent $t=2\tau+n+2$, $\abs{X_F}\prec\abs{X_P}^\lambda/(\alpha\sigma^{t-1})$,
and the smallness condition (6), the choice of $\sigma_\nu,\eta_\nu$ and the iteration estimates of \S4 are formally
unchanged; only the constant $c_0$ of (6) and hence $\gamma$ of Theorem~A change, through $b_{\rm W}(\tau)$, $C_1$ and the
factor $2$ in $\abs{X_{\hat N}}$. \textbf{The smallness exponents of Theorem~A used by W5--W6 are not affected}, and the
measure exponent of Theorem~D (W7), which depends on $\kappa$ and $\bar p-p$ only, is not affected either.
(iii) Unit conversions for W5: $\sigma=\sigma_K/(2\pi)$; for $\sigma_K\le1$ one has $\sigma\le1$, so the bound of
Lemma~\ref{bpW:w4:lattice}(iv) applies, with $\sigma^{-(2\tau+3)}=(2\pi)^{2\tau+3}\sigma_K^{-(2\tau+3)}$. The norm $\abs\cdot_r$ of
Definition~\ref{bpW:w4:taylor} differs from P96-K's in the $\theta$-component (factor $(2\pi)^{-1}$); the smallness
conditions are homogeneous in these norms, so W5 may use either consistently.
\end{remark}

\subsection{Interfaces}\label{bpW:w4:interface-sec}

\begin{remark}[From W3]\label{bpW:w4:from-w3}
After the Birkhoff map of Theorem~\ref{bpW:w3:main}, the normal coordinates are $z_j=w_{j+2}$, $\bar z_j=w_{-(j+2)}$
($j\ge1$) in the notation of Definition~\ref{bpW:w3:space}, and on real points $u_j=q_{j+2}$, $v_j=p_{j+2}$
(Lemma~\ref{bpW:w3:coords}). The restriction of the norm of $\mathfrak W$ to the normal part is
$\norm{\hat w}^2=\norm z^2+\norm{\bar z}^2=\norm u^2+\norm v^2$, with exactly the wave weights $\omega_{a,p}(j+2)$ of
Convention~\ref{bpW:w4:index}. Hence (a) no weight conversion is needed (Remark~\ref{bpW:w4:weights}); (b) the sum
norm of $D(s,r)$ and the $\mathfrak W$-norm compare by $\norm{\hat w}\le\norm u+\norm v\le\sqrt2\norm{\hat w}$; (c) the
smoothing class of W3 ($X_K$ with values in $\mathfrak W^+$) gives normal gradients in $\ell^N_{a,\bar p}$, $\bar p=p+1$, as
required by Definition~\ref{bpW:w4:taylor}; (d) the symplectic form $\Omega$ of Lemma~\ref{bpW:w3:omega} restricted to
the normal part is $\sum_jdu_j\wedge dv_j$, and $\dot z=-iH_{\bar z}$ agrees with $X_P(w)_l=-i\operatorname{sgn}(l)DP(w)e_{-l}$
(Definition~\ref{bpW:w3:hamfield-def}). The exponent $g=6$ of Remark~\ref{bpW:w3:interface} is not used in W4.

\emph{Ownership (coordinator ruling on review finding W4-R3).} W5, not W4, owns: the action--angle substitution of the
tangential modes (KP96 p.~25, Step~1), which produces $\theta$, $y$, the parameter $\xi$ and the form
$2\pi d\theta\wedge dy$; the check that the image of $D(s,r)$ lies in W3's ball $B_{\rho_1}$ (the domain check); and every
norm comparison between W3's coordinates and the KAM norms, including the sum-versus-$\mathfrak W$ comparison (b)
above. W4 keeps the shifted wave weights $\omega_{a,p}(j+2)$, so the weight-equivalence factor $3^pe^{2a}$ mentioned
in Remark~\ref{bpW:w3:interface} never arises; only the explicit phase-norm comparisons (b) remain, for W5. On this
point Remark~\ref{bpW:w3:interface}, which says ``W4 must carry these two conversions'', is superseded.
\end{remark}

\begin{remark}[To W5]\label{bpW:w4:to-w5}
W5 receives: Theorem~\ref{bpW:w4:lemma1} and Corollary~\ref{bpW:w4:lemma2} for arbitrary frequencies
$(\omega_\nu,\Omega_\nu)$ obeying \eqref{bpW:w4:freq} and $(\mathrm H_{\alpha_\nu})$ (no wave-specific input is used in
the proof); the bracket identity $R+\langle\!\langle N,F\rangle\!\rangle=\hat N+\hat a_0$ for the weak bracket
(Lemma~\ref{bpW:w4:lie}(ii)), valid for all $a,p\ge0$, and, under the guard \eqref{bpW:w4:guard}, its identification
$\langle\!\langle N,F\rangle\!\rangle=DN\cdot X_F$ with the Fr\'echet derivative of the continuous polynomial $N$
(Lemma~\ref{bpW:w4:lie}(iv)); the wave case $a=p=1$ satisfies (G), and W5 must work under (G) whenever it
differentiates $N$; and $\hat\Omega\in\ell^1_\infty$ with $\abs{\hat\Omega}_1\le\varepsilon$, so that
$\Omega_{\nu+1}=\Omega_\nu+\hat\Omega$ keeps \eqref{bpW:w4:freq} with $M_0$ increased by $\varepsilon$. W5 must supply:
the Taylor approximation $R$ of $P$ in the class $\mathcal T_s$ (P96-K (7), p.~129), including holomorphy and boundedness
of the components and of $Q^T$; the action--angle substitution, the domain check and the norm comparisons
(Remark~\ref{bpW:w4:from-w3}, ownership); the parameter regularity of $R$ in the form of the hypothesis
$\abs{X_{R(\xi)}-X_{R(\zeta)}}_r\le\varepsilon^{\mathcal L}\abs{\xi-\zeta}$; and, if it needs analyticity of the field of $P$
beyond Cauchy-type estimates, the revised W3 interface (finding W3-R1).
\end{remark}

\subsection{Ledger}\label{bpW:w4:ledger}

\paragraph{Results.} Setting: Convention~\ref{bpW:w4:index}, Definitions~\ref{bpW:w4:ops}, \ref{bpW:w4:taylor},
\ref{bpW:w4:melnikov}; wave data: Lemmas~\ref{bpW:w4:wave}, \ref{bpW:w4:k0}, Remark~\ref{bpW:w4:melnikov-check};
Fourier: Remark~\ref{bpW:w4:counterexample}, Lemmas~\ref{bpW:w4:fourier-ops}, \ref{bpW:w4:lattice}; divisors and solve:
Lemmas~\ref{bpW:w4:divisors}, \ref{bpW:w4:schur}, \ref{bpW:w4:solve}; bracket: Lemma~\ref{bpW:w4:lie}; extraction:
Lemmas~\ref{bpW:w4:extract}, \ref{bpW:w4:components}; main: Theorem~\ref{bpW:w4:lemma1}, Corollary~\ref{bpW:w4:lemma2};
comparison and interfaces: Remarks~\ref{bpW:w4:compare}, \ref{bpW:w4:from-w3}, \ref{bpW:w4:to-w5}. PLAN.md lists no
PLANNED node labels for W4; the milestone content ``operator solve, Lipschitz'' is Theorem~\ref{bpW:w4:lemma1}.

\paragraph{W1 inputs.} Lemma~\ref{bpW:w1:fourier-banach}(i),(iii) and Lemma~\ref{bpW:w1:fourier-l2}(iii) (scalar), with the
content \ref{bpW:w4:F1}--\ref{bpW:w4:F3} of Remark~\ref{bpW:w4:w1-requests}; Lemmas~\ref{bpW:w1:schur},
\ref{bpW:w1:multiplier}, \ref{bpW:w1:space}, \ref{bpW:w1:tails}; Definitions~\ref{bpW:w1:transport},
\ref{bpW:w1:kam-space}. No forward references remain.

\paragraph{Suspected source errata (P96-K).} (1) p.~127, operator-valued square-sum bound: false
(Remark~\ref{bpW:w4:counterexample}); (2) p.~127, ``$\abs{\langle k,\omega\rangle+\Omega_i-\Omega_j}\ge\alpha(1+\abs{i-j})/A_k$'': off by a
factor up to $2$ (Lemma~\ref{bpW:w4:divisors}); (3) p.~127, the (false) square-sum bound is printed with
$\abs{\!\abs{\!\abs{\tilde R}\!}\!}_{D(s)}$ unsquared on the right (typo; p.~126 has the square); (4) p.~125, $\{F,N\}+\hat N=R$ with $\hat N=[R]$
holds only modulo the constant term of $R$. Source gap: p.~128, the remaining components and the $y$-component of
$X_F$ (Theorem~\ref{bpW:w4:lemma1}).
